% concatenated from Equiv_form_rep.tex
\documentclass[11pt]{amsart}

\usepackage{geometry}
\usepackage[all]{xy}
\geometry{letterpaper}                   % ... or a4paper or a5paper or ...
%\geometry{landscape}                % Activate for for rotated page geometry
%\usepackage[parfill]{parskip}    % Activate to begin paragraphs with an empty line rather than an indent

\usepackage{epstopdf}
\usepackage{amssymb,amsmath,amsfonts,amsthm,graphicx,enumerate,amscd}
\usepackage{mathrsfs}
\usepackage{ytableau}
\usepackage{lscape}
\usepackage{hyperref}
\usepackage{rotating}

\newcommand{\tang}{\pmb{\top}}
%\newcommand{\tang}{\top}
\newcommand{\ii}{\mathrm{i}}
\newcommand{\dd}{\mathrm{d}}
\newcommand{\ee}{\mathrm{e}}
\newcommand{\smooth}{C^{\infty}}

\newcommand{\ds}{\displaystyle}
\newcommand{\deff}[1]{\textbf{\emph{\sharp1}}}

\newcommand{\func}[3]{\sharp1 \colon \sharp2 \to \sharp3}
\newcommand{\bb}[1]{\mathbb{\sharp1}}
\newcommand{\lie}[1]{\mathfrak{\sharp1}}
\newcommand{\iprod}[2]{\langle \sharp1, \sharp2 \rangle}
\newcommand{\ddell}[1]{\frac{\partial}{\partial \sharp1}}


\theoremstyle{plain}
\newtheorem{theorem}{Theorem}[section]
\newtheorem{conjecture}[theorem]{Conjecture}
\newtheorem{claim}[theorem]{Claim}
\newtheorem{corollary}[theorem]{Corollary}
\newtheorem{lemma}[theorem]{Lemma}
\newtheorem{proposition}[theorem]{Proposition}
\newtheorem*{theoremn}{Theorem}
\newtheorem*{conjecturen}{Conjecture}
\newtheorem*{claimn}{Claim}
\newtheorem*{corollaryn}{Corollary}
\newtheorem*{lemman}{Lemma}
\newtheorem*{propositionn}{Proposition}


\theoremstyle{definition}
\newtheorem{example}[theorem]{Example}
\newtheorem{remark}[theorem]{Remark}
\newtheorem{algorithm}[theorem]{Algorithm}
\newtheorem{definition}[theorem]{Definition}
\newtheorem{exercise}{}[section]
\newtheorem{problemm}{}
\newtheorem{exercisee}{Exercise}[section]
\newtheorem{fact}[theorem]{Fact}
\newtheorem{desideratum}{Desideratum}
\newtheorem*{definitionn}{Definition}
\newtheorem*{examplen}{Example}
\newtheorem*{remarkn}{Remark}
\newtheorem*{algorithmn}{Algorithm}
\newtheorem*{exercisen}{}
\newtheorem*{problemmn}{}
\newtheorem{question}[theorem]{Question}
\newtheorem*{exerciseen}{Exercise}
\newtheorem*{factn}{Fact}
\allowdisplaybreaks


\title{Equivariant formality and representation theory}
\author{Chi-Kwong Fok}
%\date{October 15, 2022}

\linespread{0.9}

\parskip     1.5mm

\begin{document}
\begin{abstract}
	%Equivariant formality, a notion in equivariant topology introduced by Goresky--Kottwitz--Macpherson, is a desirable property of spaces with group actions. Various classes of spaces of interest are well-known to be equivariantly formal, e.g., compact symplectic manifolds equipped with Hamiltonian compact Lie group actions and projective varieties equipped with linear algebraic torus actions, of which flag varieties are examples. Less is known about homogeneous spaces $G/K$ (where $G$ and $K$ are compact Lie groups) equipped with the left translation action by $K$, which is not necessarily of maximal rank. 
	Let $G$ be a compact connected Lie group and $K$ its connected Lie subgroup. Using the $K$-theoretic version of equivariant formality developed in \cite{F2} which involves analysis of vector bundles over $G/K$, we find two characterizations of equivariant formality of the isotropy action of $K$ on $G/K$. The first one equates equivariant formality with a "smoothness" condition of the restriction map of the representation ring of $G$ to that of $K$. The second characterization asserts that equivariant formality amounts to the equivariant $K$-theoretic index pairing of two certain special $K$-theory classes being nontrivial in some sense. By applying these characterizations, we are able to give a representation theoretic criterion for equivariant formality of the isotropy action by a circle subgroup, as well as an invariant theory criterion for the case where $K$ is a torus of dimension one less than the rank of $G$, culminating in a complete classification of equivariant formality where $G$ is further assumed to be simple.   %Our results unveil the $K$-theoretic nature of the symmetry of homogeneous spaces and provide a new perspective of their equivariant formality which had been studied in cohomological terms.
\end{abstract}
\maketitle
\tableofcontents
\section{Introduction}
Equivariant formality, first identified and discussed at length in \cite{BBFMP} and named in \cite{GKM}, is a desirable property of group actions on topological spaces. We say a $G$-action on a space $X$ is equivariantly formal if the Leray-Serre spectral sequence for the rational cohomology of the fiber bundle $X\hookrightarrow X\times_G EG\to BG$ associated with the Borel homotopy quotient collapses on the $E_2$ page. Interesting examples of equivariant formality abound, e.g. Hamiltonian group actions on compact symplectic manifolds and linear algebraic torus actions on smooth complex projective varieties. 

Recently, the study of equivariant formality of compact homogeneous spaces has garnered some interest. Following the convention in \cite{CFok}, we say the pair of compact Lie groups $(G, K)$ is \emph{isotropy formal} if the left action of $K$ on $G/K$ is equivariantly formal, and call such an action the \emph{isotropy action}. When $K$ is the identity subgroup, equivariant formality of isotropy action is immediate. On the other hand, the case when $K$ is of the maximal rank is well-known to be equivariantly formal. More subtle is when $K$ sits between the two extreme cases. The problem of determining isotropy formality in the intermediate cases was taken up in a series of papers. In \cite{ST, Sh}, a partial characterization of isotropy formality was given in terms of formality in the sense of rational homotopy theory and invariant theory of the equivariant cohomology coefficient rings. When $G/K$ is a generalized symmetric space, the isotropy formality of $(G, K)$ was proved in \cite{Go, GN} by comparing cohomological dimensions of the generalized symmetric space and its fixed point set on a case-by-case basis. When $K$ is a circle subgroup, \cite{Ca} provided a complete characterization of isotropy formality. In \cite{CFok}, by exploiting the notion of rational $K$-theoretic equivariant formality (RKEF) developed in \cite{F2}, the authors characterized isotropy formality in terms of the representation rings of some finite covers of $\widetilde{G}$ and $\widetilde{K}$ (such that $\pi_1(\widetilde{G})$ is torsion-free) of $G$ and $K$, and noted that the characterization in fact amounts to `smoothness' of the embedding of $\widetilde{K}$ into $\widetilde{G}$ in an algebro-geometric sense, if one takes the spec of the representation rings. In fact, considering the representation rings of some finite covers rather than the groups themselves in the characterization is for the sake of convenience as the representation ring of $\widetilde{G}$ is a (Laurent) polynomial ring and has a simpler algebraic structure in general than $R(G)$. As the property of isotropy formality is preserved under taking finite cover (\cite[Theorem 1.2]{Ca}), we expect that there should be a similar characterization where taking finite covers of the Lie groups is not necessary. In this note we will first prove this strengthened characterization where we remove the condition that $\pi_1(G)$ is torsion-free. %and offer new proofs of some results in the aforementioned papers from $K$-theoretic point of view. 
We will use $R(G; \mathbb{Q})$, $I(G; \mathbb{Q})$ and $K^*_G(X; \mathbb{Q})$ to denote the rationalized representation ring, augmentation ideal and equivariant $K$-theory respectively. %In this paper, we will present a simplified proof of a stronger representation theoretic characterization of isotropy formality (\cite[Theorem B]{CFok}) where we remove the condition that $\pi_1(G)$ be torsion-free.
\begin{theorem}\label{rep}
	Let $G$ be a compact connected Lie group and $K$ its connected Lie subgroup. Then $(G, K)$ is isotropy formal if and only if the image $R:=i^*R(G; \mathbb{Q})$ of the restriction map $i^*: R(G;\mathbb{Q})\to R(K;\mathbb{Q})$ is regular at the restriction  $I:=i^*I(G;\mathbb{Q})$ of the augmentation ideal of $R(G;\mathbb{Q})$.
\end{theorem}
As remarked earlier and in \cite[Introduction]{CFok}, the isotropy formality of $(G, K)$, through the regularity condition in the above characterization, can be recast as the smoothness at the point $I(G; \mathbb{Q})$ of the map $\text{Spec}R(K; \mathbb{Q})\to\text{Spec}R(G; \mathbb{Q})$, which is the affinization of the embedding $K\hookrightarrow G$. As a corollary, the rational representation ring $R(G;\mathbb{Q})$ is always regular at the augmentation ideal, since $(G, G)$ is obviously isotropy formal. In particular, taking the quotient of a compact connected Lie group $G$ by its finite central subgroup $Z$ does not create a `crease' at $I(G/Z; \mathbb{Q})$ in $\text{Spec}R(G/Z; \mathbb{Q})$. It was claimed after \cite[Proof of Lemma 4.2]{Ho} that $R(\text{PSU}(3); \mathbb{Q})$ is not regular at its augmentation ideal, citing the existence of torsion in $\pi_1(\text{PSU}(3))$. We will show by explicit computation that $R(\text{PSU}(3); \mathbb{Q})$ is in fact regular at its augmentation ideal, which is consistent with Theorem \ref{rep}. This will serve as a correction to the claim after \cite[Lemma 4.2]{Ho}.

%In proving Theorem \ref{rep}, we obtain a description of what we call the \emph{equivariant $K$-theoretic Samelson subspace}, which roughly speaking is the primitive subspace of the exterior algebra part of the $K$-theory ring $K_K^*(G/K; \mathbb{Q})$ (Lemma \ref{eqSam}). This is used to describe the ordinary $K$-theoretic Samelson subspace of $K^*(G/K)$ for $(G, K)$ isotropy formal. Another ingredient in the proof of Theorem \ref{rep} is \cite[Theorem A]{CFok}, which is inspired by the results in \cite{ST, Sh} and links isotropy formality with formality of the homogeneous space\footnote{One shall note that equivariant formality does not imply formality in general: consider the trivial action on a non-formal space by a Lie group.}
%\begin{theorem}(\cite[Theorem A]{CFok})\label{formality}
%	Assume the same conditions as in Theorem \ref{rep}. If $(G, K)$ is isotropy formal, then $G/K$ is formal.
%\end{theorem}
%Whereas the proof in \cite{CFok} invoked a pure Sullivan model for the equivariant and ordinary cohomology rings of $G/K$ and the complete intersection condition of a certain subring of $H^*(G/K; \mathbb{Q})$ to yield formality of $G/K$, ours employs a technical lemma about the structure of the equivariant $K$-theory of equivariantly formal spaces to show another equivalent condition for formality of $G/K$, i.e. the Samelson subspace reaching the maximal possible dimension (see Proposition \ref{kthrystructure}). This result, together with our description of (equivariant) $K$-theoretic Samelson subspace and the characterization of equivariant formality using RKEF, culminates in the proof of Theorem \ref{rep}. 

To demonstrate the utility of Theorem \ref{rep}, we apply it to obtain characterizations of isotropy formality in two contexts.  The first one is when $K$ is a circle subgroup of $G$, which is assumed to be semisimple. We characterize isotropy formality for this case in Theorem \ref{circle}. Interestingly, Theorem \ref{rep} reduces the proof of Theorem \ref{circle} to elementary calculus of curves when checking the smoothness of the map $\text{Spec}R(K; \mathbb{Q})\to\text{Spec}R(G; \mathbb{Q})$. In fact this result recovers \cite[Theorem 1.5]{Ca} (in the case where $S$ is contained the commutator subgroup $G'$, i.e., the semisimple part of $G$), where the characterization is in terms of the ``reflectibility'' of $S$ and whose proof involves cohomological arguments.
\begin{theorem}\label{circle}
	Let $S$ be a circular subgroup of a compact connected semisimple Lie group $G$. Then $(G, S)$ is isotropy formal if and only if any representation of $G$, when restricted to $S$, is self-dual. 
\end{theorem}
We also apply this Theorem \ref{rep} to the case where $K$ is a torus subgroup of dimension being 1 less than the rank of $G$. Via the Chern character, we translate the regularity criterion for representation rings to the regularity criterion below for a distinguished polynomial on the Lie algebra of $K$ invariant under the Weyl group, which is a computable algorithm for determining whether an any given corank 1 pair is isotropy formal.

\begin{theorem}[=Theorem \ref{invtpoly}]
	Let $G$ be a compact connected Lie group of rank $n$, $S$ a torus subgroup of dimension being $n-1$, $T$ a maximal torus of $G$ containing $S$, $\mu\in \chi(T)$ a primitive character of $T$ such that $S=\text{ker }\mu$, $L$ a primitive weight in the weight lattice $\mathfrak{t}^*_\mathbb{Z}$ corresponding to $\mu$ after exponentiation, and $W$ the Weyl group of $G$. Define
	\[\Phi_L:=\left(\prod_{w\in W/W_{[L]}} wL\right)^{n(L)}\in \text{Sym}^*(\mathfrak{t}^*)^W,\]
	where $W_{[L]}$ is the stabilizer of $[L]:=\{\pm L\}\in \mathfrak{t}^*/\{\pm 1\}$, $n(L)=1$ or 2 depending on whether $\displaystyle\prod_{w\in W/W_{[L]}}wL$ is $W$-invariant or $W$-anti-invariant, i.e. 
	\[\displaystyle w'(\prod_{w\in W/W_{[L]}}wL)=\text{sgn}(w')\prod_{w\in W/W_{[L]}}wL\text{ for all }w'\in W.\] 
	Then $(G, S)$ is isotropy formal if and only if $\Phi_L\notin (f_1, \cdots, f_n)^2$, where $f_1, \cdots, f_n$ are the fundamental $W$-invariant homogeneous polynomials on $\mathfrak{t}$, i.e. $\text{Sym}^*(\mathfrak{t}^*)^W=\mathbb{R}[f_1, \cdots, f_n]$.
\end{theorem}
Theorem \ref{invtpoly}, together with a sharp bound on the degree on the distinguished polynomial (Lemma \ref{degreebound}), leads us to the following classification of isotropy formal action of corank 1.
\begin{theorem}\label{corank1}
		Let $G$ be a compact and connected simple Lie group of rank $n$, $S$ a torus subgroup of dimension being $n-1$, $T$ a maximal torus of $G$ containing $S$, $\mu\in \chi(T)$ a primitive character of $T$ such that $S=\text{ker }\mu$, $L$ a primitive weight in the weight lattice $\mathfrak{t}^*_\mathbb{Z}$ corresponding to $\mu$ after exponentiation, and $W$ the Weyl group of $G$. The following table classifies the isotropy formality of $(G, S)$ in terms of $L$. Here $\omega_i$ is the $i$-th fundamental weight of the relevant simple Lie group following the Bourbaki convention.
		\begin{center}
		\begin{tabular}{|c|c|}
		\hline
		$G$&$L$\\
		\hline
		$A_n, n\geq 2, n\neq 3$& $L\in W\cdot \omega_1$\\
		\hline
		$A_3$&$L\in W\cdot \omega_1\cup W\cdot \omega_2$\\
		\hline
		$B_2, C_2$&$\text{Any primitive weight}$\\
		\hline
		$B_n, C_n, n=3, 4$&$L\in W\cdot \omega_1\cup W\cdot \omega_n$\\
		\hline
		$B_n, C_n, n\geq 5$&$L\in W\cdot \omega_1$\\
		\hline
		$D_4$&$L\in W\cdot\omega_1\cup W\cdot \omega_3\cup W\cdot \omega_4$\\
		\hline
		$D_n, n\geq 5$&$L\in W\cdot \omega_1$\\
		\hline
		$G_2$&$\text{Any primitive weight}$\\
		\hline
		$F_4, E_6, E_7, E_8$&$\text{No such primitive weights}$\\
		\hline
		\end{tabular}
		\end{center}
\end{theorem}
It should be noted that a characterization of isotropy formality of general corank 1 pairs $(G, K)$ is given in \cite{CH}. This characterization consists of elaborate steps of checking the fundamental groups and total Betti number of a certain homogeneous space. The authors approach this characterization by way of a classification of $G/K$ with the rational homotopy type of a product of an odd- and an even-dimensional sphere. Table 1 in \cite{CH}, which is part of the characterization algorithm, determines isotropy formality for  irreducible compact connected pairs $(G, K)$ with $G/K$ simply-connected and its rational cohomology isomorphic to that of the product of two spheres. The isotropy formal pairs in the table are consistent with those listed in Theorem \ref{corank1}. %However, we, while Table 1 in \cite{CH} does not contain cases 

We also find another characterization of isotropy formality with a different flavor. We first observe that formality of $G/S$ (where $S$ is a torus subgroup) is equivalent to the nontrivial $K$-theoretic pairing of certain two special $K$-theory classes (Proposition \ref{nondegenformal}). Based on this we obtain the following equivariant analogue. 
\begin{theorem}\label{indexthry}
	Let $G$ be a compact connected Lie group and $S$ a torus subgroup, and $\pi: G/S\to\text{pt}$ be the collapsing map. Denote by $A$ the subring of $K_S^*(G/S)\otimes\mathbb{Q}$ generated by $S$-equivariant line bundles $G\times_S\mathbb{C}_\mu$ where $\mu\in\mathfrak{s}^*$. Define $\mathcal{L}_S$ to be the $R(S; \mathbb{Q})$-submodule of $K_S^*(G/S)\otimes\mathbb{Q}$ generated by the odd equivariant $K$-theory classes of $G/S$ defined by the $\delta$-construction (see Definitions \ref{deltaconstruction} and \ref{equivsamelson}). Then $(G, S)$ is isotropy formal if and only if there exist $a\in A$ and $b\in\bigwedge\nolimits_{R(S; \mathbb{Q})}^{\text{rank }G-\text{rank }S}\mathcal{L}_S$ such that $\pi_*(ab)\notin I(S; \mathbb{Q})$. 
\end{theorem}
We end this paper by discussing the isotropy formality of two examples of pairs using the aforementioned main results.

Throughout this paper, $G$ always means a compact connected Lie group and $K$ its connected Lie subgroup unless otherwise specified. 

\textbf{Acknowledgment }We would like to gratefully acknowledge the support by the Xi'an Jiaotong-Liverpool University Research Development Fund RDF 23-02-029.
%Throughout this paper, all (equivariant) $K$-theory $K_G^*(X)$ and representation rings $R(G)$ should be understood to be rationalized, i.e., $K_G^*(X)\otimes\mathbb{Q}$ and $R(G)\otimes\mathbb{Q}$ respectively, unless otherwise specified.
%It should be noted that the other case covered by \cite[Theorem 1.5]{Ca} where $S$ is not contained in $S$ can be shown using 
\section{Preliminaries}
In this section, we will recall a number of definitions and results about the (equivariant) $K$-theory of compact homogeneous spaces and $K$-theoretic equivariant formality, which are all crucial ingredients in the proofs of the main results of this paper.
\begin{definition}(\cite[\S I.4]{Ho}, \cite[\S 3]{BZ}, \cite[Definitions 2.2 and 2.5]{F})\label{deltaconstruction}
	Let $\delta: R(G)\to K^{-1}(G)$ be the $\mathbb{Z}$-linear map which sends a complex representation $\rho$ with underlying vector space $V$ to the $K$-theory class represented by the complex of vector bundles
	\begin{align*}
		0\longrightarrow G\times\mathbb{R}\times V&\longrightarrow G\times\mathbb{R}\times V\longrightarrow 0\\
		(g, t, v)&\mapsto\begin{cases}(g, t, -t\rho(g)v),&\ \text{if }t>0\\ (g, t, tv),&\ \text{if }t\leq 0.\end{cases}
	\end{align*}
	More generally, for an element $\rho\in\text{ker}(i^*: R(G)\to R(K))$ which is the formal difference $\rho_1-\rho_2\in R(G)$ of two complex $G$-representations restricting to the same $K$-representations, we define the $\mathbb{Z}$-linear map $\delta^{G/K}:\text{ker}(i^*)\to K^{-1}(G/K)$ which sends $\rho$ to the complex of vector bundles
	\begin{align*}
		0\longrightarrow G/K\times\mathbb{R}\times V&\longrightarrow G/K\times\mathbb{R}\times V\longrightarrow 0\\
		(gK, t, v)&\mapsto\begin{cases}(gK, t, t\rho_1(g)\rho_2(g)^{-1}v),&\ \text{if }t>0\\ (gK, t, tv),&\ \text{if }t\leq 0.\end{cases}
	\end{align*}
	The map $\delta_K^{G/K}: \text{ker}(i^*)\to K_K^{-1}(G/K)$ is defined using the same complex of vector bundles and equipping it with the $K$-action on $G/K\times \mathbb{R}\times V$ given by $k\cdot(gK, t, v)=(kgK, t, \rho_1(k)v)$. The map $\delta_K^G: R(G)\to K_K^{-1}(G)$ can be defined similarly (here $K$ acts on $G$ by conjugation). We will also use the same notation to denote the rationalized maps from the rational representation ring to the rational $K$-theory $K^{-1}(G; \mathbb{Q})$, $K^{-1}(G/K; \mathbb{Q})$ etc.
\end{definition}
\begin{definition}(\cite[p. 172]{H}, \cite{GD}). Let $A$ be a ring and $B$ an $A$-algebra. The module of K\"ahler differentials $\Omega^1_{B/A}$ of $B$ over $A$ is the quotient of the $B$-module freely generated by the symbols $\{db| b\in B\}$ by the submodule generated by the following sets
	\begin{enumerate}
		\item(Constancy of elements of $A$) $\{da| a\in A\}$, 
		\item(Linearity of $d$) $\{d(b_1+b_2)-db_1-db_2| b_1, b_2\in B\}$, and 
		\item(Derivation) $\{d(b_1b_2)-b_1db_2-b_2db_1| b_1, b_2\in B\}$.
	\end{enumerate}
	The Grothendieck differentials $\Omega^*_{B/A}$ is defined to be the exterior $B$-algebra $\displaystyle\bigoplus_{p=0}^\infty \bigwedge\nolimits_B^p\Omega_{B/A}^1$.
\end{definition}
\begin{theorem}\label{kthrygrp}
	\begin{enumerate}
		\item \textnormal{(\cite[Theorem 2.1]{Ho}).}\label{Ho} The map $\delta$ satisfies
		\[\delta(\rho_1\otimes\rho_2)=\text{dim}(\rho_1)\delta(\rho_2)+\text{dim}(\rho_2)\delta(\rho_1).\]
		Moreover, $\delta$ induces an isomorphism of Hopf algebras
		\[\bigwedge\nolimits_\mathbb{Q}^*(\text{Im}(\delta)\otimes\mathbb{Q})\cong K^*(G; \mathbb{Q}).\]
		\item \textnormal{(cf. \cite{BZ})}\label{BZ} The map 
		\begin{align*}
			\Omega^*_{R(G)/\mathbb{Z}}\otimes_{R(G)}R(K)&\to K_K^*(G)\\
			d\rho_1\otimes\rho_2&\mapsto \rho_2\cdot\delta_K^G(\rho_1)
		\end{align*}
		is injective. In particular, when $\pi_1(G)$ is torsion-free, the map is an isomorphism.
	\end{enumerate}
\end{theorem}
\begin{remark}\label{deltaredux}
	We have that $\Omega^1_{R(G; \mathbb{Q})/\mathbb{Q}}\otimes_{R(G; \mathbb{Q})}\mathbb{Q}\cong\text{Im}(\delta)\otimes\mathbb{Q}$ via the map $d\rho\otimes 1\mapsto \delta(\rho)$, and that $\delta^{G/K}$ also satisfies the derivation property of $\delta$ as in Theorem \ref{kthrygrp} (\ref{Ho}).
\end{remark}
\begin{proposition}\label{log}
	Let $\rho_1\cdot\rho_2: G\to \text{GL}(V)$ be the pointwise multiplication $(\rho_1\cdot\rho_2)(g):=\rho_1(g)\rho_2(g)$. Then $\delta(\rho_1\cdot\rho_2)=\delta(\rho_1)+\delta(\rho_2)$.
\end{proposition}
\begin{proof}
	This is an immediate consequence of the fact that the two maps $G \to \text{GL}(V\oplus V)$ defined by 
	\[g\mapsto \begin{pmatrix}\rho_1(g)& 0\\ 0&\rho_2(g)\end{pmatrix}\ \text{and }g\mapsto \begin{pmatrix}(\rho_1\cdot\rho_2)(g)&0\\ 0& \text{I}_V\end{pmatrix}\]
	are $G$-equivariantly homotopic (cf. \cite[Proof of Proposition 3.1]{BZ}) and that $\delta(n)=0$ for $n\in\mathbb{Z}$.
\end{proof}
\begin{proposition}\label{repringreg}
	$\text{dim}_\mathbb{Q}\Omega^1_{R(G; \mathbb{Q})/\mathbb{Q}}\otimes_{R(G; \mathbb{Q})}\mathbb{Q}=\text{rank }G$
\end{proposition}
\begin{proof}
	The LHS is isomorphic to $\text{Im}(\delta)\otimes\mathbb{Q}$, which is the primitive space of $K^*(G; \mathbb{Q})$ by Theorem \ref{kthrygrp} (\ref{Ho}). Note that $\text{dim }K^*(G; \mathbb{Q})=\text{dim }H^*(G; \mathbb{Q})=2^{\text{rank }G}$ (for the second equality, see, for example, \cite[Theorem 3.6]{F3}), so the dimension of the primitive space is $\text{rank }G$. 
\end{proof}
%\begin{lemma}
	%Let $k$ be a field of characteristic 0, $A$ a finitely generated $k$-algebra and $f: A\to B$ a $k$-algebra homomorphism. Then $\text{rank}_A\Omega^1_{A/k}\leq \text{rank}_B\Omega^1_{A/k}\otimes_A B$. 
%\end{lemma}
\begin{proposition}\label{rkkahler}
	$\text{rank}_{R(K; \mathbb{Q})}\Omega^1_{R(G; \mathbb{Q})/\mathbb{Q}}\otimes_{R(G; \mathbb{Q})}R(K; \mathbb{Q})=\text{rank }G$.
\end{proposition}
\begin{proof}
	Let $R:=i^*R(G; \mathbb{Q})$. It suffices to show that $\text{rank}_R\Omega^1_{R(G; \mathbb{Q})/\mathbb{Q}}\otimes_{R(G; \mathbb{Q})}R=\text{rank}G$ because $R$ is a subring of $R(K; \mathbb{Q})$. As $R(K; \mathbb{Q})$ is a subring of the (rationalized) representation ring of a maximal torus of $K$, which is an integral domain, $R$ is also an integral domain. Then $\text{ker}(i^*)$ is a prime ideal of $R(G; \mathbb{Q})$. Moreover, $\text{ker}(i^*)\subseteq I(G; \mathbb{Q})$. For simplicity we denote $\text{ker}(i^*)$ by $\mathfrak{p}$. By \cite[Lemmas 7.12(2) and 7.13]{CFok}, we have
	\[\text{rank}_{R(G; \mathbb{Q})}\Omega^1_{R(G; \mathbb{Q})/\mathbb{Q}}=\text{rank}_{R(G; \mathbb{Q})_{\mathfrak{p}}}\Omega^1_{R(G; \mathbb{Q})/\mathbb{Q}}\otimes_{R(G; \mathbb{Q})}R(G; \mathbb{Q})_{\mathfrak{p}}=\text{rank}G.\]
	It follows that 
	\begin{align*}
		&\text{rank}G\\
		=&\text{rank}_{R(G; \mathbb{Q})_{\mathfrak{p}}}\Omega^1_{R(G; \mathbb{Q})/\mathbb{Q}}\otimes_{R(G; \mathbb{Q})}R(G; \mathbb{Q})_{\mathfrak{p}}\\
		\leq &\text{rank}_{R(G; \mathbb{Q})/\mathfrak{p}}\Omega^1_{R(G; \mathbb{Q})/\mathbb{Q}}\otimes_{R(G; \mathbb{Q})}(R(G; \mathbb{Q})_{\mathfrak{p}}/\mathfrak{p}R(G; \mathbb{Q})_{\mathfrak{p}})\ (\text{By Nakayama's Lemma})\\
		=&\text{rank}_{R(G; \mathbb{Q})/\mathfrak{p}}\Omega^1_{R(G; \mathbb{Q})/\mathbb{Q}}\otimes_{R(G; \mathbb{Q})}(R(G; \mathbb{Q})_{I(G; \mathbb{Q})}/\mathfrak{p}R(G; \mathbb{Q})_{I(G; \mathbb{Q})})\\
		&(\text{As }R(G; \mathbb{Q})_{\mathfrak{p}}/\mathfrak{p}R(G; \mathbb{Q})_{\mathfrak{p}}\cong R(G; \mathbb{Q})_{I(G; \mathbb{Q})}/\mathfrak{p}R(G; \mathbb{Q})_{I(G; \mathbb{Q})})\\
		\leq&\text{rank}_{R(G; \mathbb{Q})/I(G; \mathbb{Q})}\Omega^1_{R(G; \mathbb{Q})/\mathbb{Q}}\otimes_{R(G; \mathbb{Q})}(R(G; \mathbb{Q})_{I(G; \mathbb{Q})}/I(G; \mathbb{Q})R(G; \mathbb{Q})_{I(G; \mathbb{Q})})\\
		&(\text{By Nakayama's Lemma})\\
		=&\text{rank}_{\mathbb{Q}}\Omega^1_{R(G; \mathbb{Q})/\mathbb{Q}}\otimes_{R(G; \mathbb{Q})}\mathbb{Q}\\
		=&\text{rank}G\ (\text{By Proposition \ref{repringreg}}).
	\end{align*}
	Thus 
	\begin{align*}
		\text{rank}_R\Omega^1_{R(G; \mathbb{Q})/\mathbb{Q}}\otimes_{R(G; \mathbb{Q})}R&=\text{rank}_{R(G; \mathbb{Q})/\mathfrak{p}}\Omega^1_{R(G; \mathbb{Q})/\mathbb{Q}}\otimes_{R(G; \mathbb{Q})}(R(G; \mathbb{Q})_{\mathfrak{p}}/\mathfrak{p}R(G; \mathbb{Q})_{\mathfrak{p}})\\
		&=\text{rank}G
	\end{align*}
	as desired.
\end{proof}
\begin{definition}(cf. \cite[Definition 6.1]{CFok}).\label{equivsamelson}
	Let $Q$ be the $K$-theoretic Samelson space, i.e., the space of primitive elements of $K^*(G; \mathbb{Q})$ in the image of the pullback map $j^*: K^*(G/K; \mathbb{Q})\to K^*(G; \mathbb{Q})$ induced by the quotient map $j: G\to G/K$. Similarly, let $Q_K$ be the equivariant $K$-theoretic Samelson space, i.e., $\text{Im}(j_K^*)\cap \mathcal{N}_K$, where $j_K^*: K_K^*(G/K; \mathbb{Q})\to K_K^*(G; \mathbb{Q})$ is the pullback map on equivariant $K$-theory, and $\mathcal{N}_K$ the $\mathbb{Q}$-vector subspace of $K_K^*(G; \mathbb{Q})$ spanned by the primitive elements $\{\delta_K^G(\rho)|\rho\in R(G; \mathbb{Q})\}$. Let $\mathcal{M}_K$ be the $\mathbb{Q}$-vector subspace of $K_K^*(G; \mathbb{Q})$ generated by $\{\delta_K^G(\rho)|\rho\in\text{ker}(i^*)\}$, and $\mathcal{L}_K$ the $R(K; \mathbb{Q})$-submodule of $K_K^*(G/K; \mathbb{Q})$ generated by $\{\delta_K^{G/K}(\rho)|\rho\in\text{ker}(i^*)\}$.%Define $P$ to be the preimage $(j^*)^{-1}(Q)$. 
\end{definition}
The following result about the structure of the $K$-theory of $G/K$ is deduced from the structure of its cohomology (\cite[pp. 73, 83, 152]{GHV}) and by applying the Chern character isomorphism.
\begin{proposition}\textnormal{(cf. \cite[Theorem 6.2]{CFok}).}\label{kthrystructure}
	We have the ring isomorphism
	\[K^*(G/K; \mathbb{Q})\cong (R(K; \mathbb{Q})/i^*I(G; \mathbb{Q})\oplus\mathfrak{a})\otimes\bigwedge\nolimits^*P.\]
	Here $R(K; \mathbb{Q})/i^*(G; \mathbb{Q})$ is the image of the map $\alpha: R(K; \mathbb{Q})\to K^0(G/K; \mathbb{Q})$ defined by the associated vector bundle $\alpha(\rho):=[G\times_K V]$, $\mathfrak{a}$ is an ideal of the first tensor factor on the RHS. We also have that $j^*$ maps $P$ isomorphically onto $Q$. Moreover, $\text{dim }P=\text{dim }Q\leq \text{rank }G-\text{rank }K$, and equality holds if and only if $G/K$ is formal, if and only if $\mathfrak{a}=0$.
\end{proposition}
\begin{lemma}\label{key}
	We have $\delta^{G/K}(\text{ker}(i^*))\subseteq P$.% and $Q\subseteq\{\delta(\rho)|\delta(i^*(\rho))=0\in K^{-1}(K; \mathbb{Q})\}$.
\end{lemma}
\begin{proof}
	Let $x\in\text{ker}(i^*)$ be the formal difference $\rho_1-\rho_2$ of complex $G$-representations restricting to the same $K$-representation. Then $j^*(\delta^{G/K}(x))=\delta(\rho_1\cdot\rho_2^{-1})=\delta(\rho_1)-\delta(\rho_2)\in Q$ (the last equality holds by Proposition \ref{log}) and thus $\delta^{G/K}(x)\in P$. It follows that $\delta^{G/K}(\text{ker}(i^*))\subseteq P$.
%For the second assertion, we let $p\in P$ and consider an arbitrary (virtual) representation $\rho:=\rho_1-\rho_2\in\delta^{-1}(j^*(p))\subseteq R(G; \mathbb{Q})$, where $\rho_1$ and $\rho_2$ are positive rational linear combinations of irreducible representations of $G$. Then there exists a positive integer $m$ such that $\sigma_1:=m\rho_1$ and $\sigma_2=m\rho_2$ are actual $G$-representations. Let $i: K\to G$ and $i': \text{pt}\to G/K$ be inclusion maps (the latter being inclusion as $eK$) and $j': K\to \text{pt}$ be the collapsing map. Then $j\circ i=i'\circ j'$ and 
	%\[\delta(i^*(\sigma_1-\sigma_2))=i^*\delta(\sigma_1-\sigma_2)=i^*(j^*(mp))=j'^*(i'^*(mp))=j'^*(0)=0.\]
	%Now the second assertion follows.%Since $\sigma_1$ and $\sigma_2$ have the same dimension, we conclude that $i^*(\sigma_1-\sigma_2)=0$, $\rho_1-\rho_2\in\text{ker}(i^*)$ and establish the reverse inclusion $P\subseteq\delta^{G/K}(\text{ker}(i^*))$. This finishes the proof of the proposition.
\end{proof}
A main ingredient in our representation theoretic proof of Theorem \ref{rep} is the following $K$-theoretic characterization of equivariant formality.
\begin{theorem}\textnormal{(\cite[Theorem 1.3]{F2}).}\label{RKEF} Let $G$ act on a finite CW-complex $X$. Then the following are equivalent. 
\begin{enumerate}
	\item $X$ is an equivariantly formal $G$-space.
	\item The forgetful map $f: K_G^*(X; \mathbb{Q})\to K^*(X; \mathbb{Q})$ is surjective. 
	\item The forgetful map $f$ induces an isomorphism $K_G^*(X; \mathbb{Q})\otimes_{R(G; \mathbb{Q})}\mathbb{Q}\to K^*(X; \mathbb{Q})$.
\end{enumerate}
\end{theorem}
Later on we will, by abuse of notation, always use $f$ to denote the forgetful map for any topological spaces $X$. 
\section{Proof of Theorems \ref{rep} and \ref{circle}} 
\begin{lemma}\label{eqSam}
	The equivariant $K$-theoretic Samelson subspace $Q_K$ is the $\mathbb{Q}$-vector subspace $\mathcal{M}_K$ of $K_K^*(G; \mathbb{Q})$.
\end{lemma}
\begin{proof}
	Let us first prove the case where $K$ is a ($k$-dimensional) torus subgroup $S$. The vector subspace $\mathcal{M}_S$ is in $Q_S$ because $j_S^*(\delta_K^{G/K}(\rho))=\delta_K^G(\rho)$ if $\rho\in\text{ker}(i^*)$. The remaining task is to show that $Q_S\subseteq\mathcal{M}_S$. Let $r_S:S\to G$ and $r_{S/S}: S/S\to G/S$ be inclusion maps, and $k_S:S\to S/S$ be the quotient map. Pick an arbitrary element in $Q_S$. Multiplying an integer if necessary, we may assume that the element is of the form $\delta_S^G(\rho_1-\rho_2)$ where $\rho_1$ and $\rho_2$ are representations of $G$. Let $x\in K^{-1}_S(G/S; \mathbb{Q})$ be a preimage of $\delta_S^G(\rho_1-\rho_2)$ under $j_S^*$. Then we have 
	\begin{align*}
		r_S^*(\delta_S^G(\rho_1-\rho_2))&=r_S^*(j_S^*(x))\\
								&=k_S^*(r_{S/S}^*(x))\ (\text{because }j_S\circ r_S=r_{S/S}\circ k_S)\\
								&=0\ (\text{because }r_{S/S}^*(x)\in K_S^{-1}(S/S; \mathbb{Q})=0).
	\end{align*}
	On the other hand, suppose $\rho_1$ and $\rho_2$ decomposes into the direct sums of irreducible 1-dimensional representations of $S$, $\displaystyle\alpha:=\sum_m t_1^{\alpha_{1m}}t_2^{\alpha_{2m}}\cdots t_k^{\alpha_{km}}$ and $\displaystyle \beta:=\sum_n t_1^{\beta_{1n}}t_2^{\beta_{2n}}\cdots t_k^{\beta_{kn}}$ on restriction to $S$. Then by a generalization of the description of the restriction map $K_T^*(G)\to K_T^*(T)$ with $T$ being a maximal torus of $G$ (cf. \cite[Lemma 4.19]{F}), we have
	\begin{align*}
		&r_S^*(\delta_S^G(\rho_1-\rho_2))\\
		=&\sum_m t_1^{\alpha_{1m}}t_2^{\alpha_{2m}}\cdots t_k^{\alpha_{km}}\otimes\left(\sum_{i=1}^k\alpha_{im}\delta(t_i)\right)-\sum_n t_1^{\beta_{1n}}t_2^{\beta_{2n}}\cdots t_k^{\beta_{kn}}\otimes\left(\sum_{i=1}^n\beta_{in}\delta(t_i)\right)\\
		=&\sum_{i=1}^k t_i\left(\frac{\partial \alpha}{\partial t_i}-\frac{\partial\beta}{\partial t_i}\right)\otimes\delta(t_i)
	\end{align*}
	Thus $\displaystyle\frac{\partial}{\partial t_i}(\alpha-\beta)=0$ for all $1\leq i\leq k$, and $\alpha$ and $\beta$ differ by a trivial representation of dimension say, $\ell$. Then $\delta_S^G(\rho_1-\rho_2)=\delta_S^G(\rho_1-\rho_2-\ell)$ and $\rho_1-\rho_2-\ell\in\text{ker}(i^*)$. This shows $Q_S\subseteq \mathcal{M}_S$.
	
	%Let $Z_G(S)_0$ be the connected component of the centralizer of $S$ in $G$. Then $Z_G(S)_0/S$ is a Lie group. Let $\widetilde{Z}$ be the connected Lie subgroup of $Z_G(S)_0$ which has the Lie algebra isomorphic to that of $Z_G(S)_0/S$. Then $\widetilde{Z}$ is a finite covering group of $Z_G(S)_0/S$ and so is $\widetilde{Z}\times S$ of $Z_G(S)_0$. Consider the following commutative diagram. 
	%\begin{eqnarray}
	%	\xymatrix{K_S^*(G/S; \mathbb{Q})\ar[r]^{j_S^*}\ar[d]^{r_1^*}&K_S^*(G; \mathbb{Q})\ar[d]^{r_2^*}\\ K_S^*(Z_G(S)_0/S; \mathbb{Q})\ar[r]^{k^*}\ar[d]^{p_1^*}& K_S^*(Z_G(S)_0; \mathbb{Q})\ar[d]^{p_2^*}\\ K_S^*(\widetilde{Z}; \mathbb{Q})\ar[r]^{\ell^*}\ar[d]^\cong& K_S^*(\widetilde{Z}\times S; \mathbb{Q})\ar[d]^{\cong}\\ R(S; \mathbb{Q})\otimes K^*(\widetilde{Z}; \mathbb{Q})\ar[r]& R(S; \mathbb{Q})\otimes K^*(S; \mathbb{Q})\otimes K^*(\widetilde{Z}; \mathbb{Q})}
	%\end{eqnarray}
	%Here $p_1^*$, $p_2^*$ and $\ell^*$ are maps induced by projection and $r_1^*$, $r_2^*$ and $k^*$ are restriction maps. Let $\displaystyle\sum_i\mu_i\delta_S^G(\rho_i)\in Q_S$, where $\mu_i\in R(S; \mathbb{Q})$ and $\rho_i$ is a representation of $G$. Suppose $\rho_i|_{Z_G(S)_0}$ is decomposed into irreducible representations of $Z_G(S)_0$ as $\displaystyle\sum_j \rho_{ij}$. The map $p_2^*: R(Z_G(S)_0; \mathbb{Q})\to R(\widetilde{Z}\times S; \mathbb{Q})\cong R(\widetilde{Z}; \mathbb{Q})\otimes R(S; \mathbb{Q})$ induced by the projection of the covering map is injective, and sends $\rho_{ij}$ to the tensor product $\rho'_{ij}\otimes \mu_{ij}$ of its restrictions to $\widetilde{Z}$ and $S$ which are irreducible as well. By the derivation property of the $\delta$-map, a generalization of the description of the restriction map $K_T^*(G)\to K_T^*(T)$ with $T$ being a maximal torus of $G$ (cf. \cite[Lemma 4.19]{F}) and identifying $K_S^*(\widetilde{Z}\times S; \mathbb{Q})$ with $R(S; \mathbb{Q})\otimes K^*(S; \mathbb{Q})\otimes K^*(\widetilde{Z}; \mathbb{Q})$ (cf. \cite[Proposition 2.2]{Se2}), we have
	%\[p_2^*(r_2^*(\delta_S^G(\rho_i)))=\sum_j(\mu_{ij}\otimes\delta(\mu_{ij})\otimes\text{dim }\rho'_{ij}+1\otimes 1\otimes \delta(\rho'_{ij})).\]
	%By the commutativity of the diagram, $p_2^*(r_2^*(\delta_S^G(\rho_i)))\in p_2^*(r_2^*(Q_S))\subseteq\text{Im}(\ell^*)$. With the identifications in the bottom left and right of the diagram (cf. \cite[Proposition 2.2]{Se2}), we have that $\text{Im}(\ell^*)=R(S; \mathbb{Q})\otimes 1\otimes K^*(\widetilde{Z}; \mathbb{Q})$. It follows that $\delta(\mu_{ij})=0\in K^*(S; \mathbb{Q})$ and hence $\mu_{ij}$ is the trivial one-dimensional representation of $S$. In other words, $i^*(\rho_{ij})$ are trivial representations and so are $i^*(\rho_i)$ for all $i$ and $j$. Note that $\displaystyle\sum_i\mu_i\delta_S^{G/S}(\rho_i)=\sum_i\mu_i\delta_S^{G/S}(\rho_i-{\text{dim}\rho_i})$ and thus it is in $\mathcal{M}_S$. We then have shown that $\mathcal{M}_S=Q_S$. 
	
	As to the general case, we let $S$ be a maximal torus of $K$ and consider the following commutative diagram.
	\begin{eqnarray}
		\xymatrix{K_K^*(G/K; \mathbb{Q})\ar[r]^{j_K^*}\ar[d]_{f_1^*}& K_K^*(G; \mathbb{Q})\ar[d]^{f_2^*}\\ K_S^*(G/S; \mathbb{Q})\ar[r]^{j_S^*}&K_S^*(G; \mathbb{Q})}
	\end{eqnarray}
	Here $f_1^*$ is the composition of the map $K_K^*(G/K; \mathbb{Q})\to K_S^*(G/K; \mathbb{Q})$ induced by $R(K; \mathbb{Q})\to R(S; \mathbb{Q})$ and the map $K_S^*(G/K; \mathbb{Q})\to K_S^*(G/S; \mathbb{Q})$ induced by the projection $G/S\to G/K$, while $f_2^*$ is induced by $R(K; \mathbb{Q})\to R(S; \mathbb{Q})$. Again it is enough to consider $\delta_K^G(\rho_1-\rho_2)\in Q_K$, where $\rho_1$ and $\rho_2$ are $G$-representations, and show that $\rho_1|_K$ and $\rho_2|_K$ differ by a trivial representation. Note that $f_2^*(\delta_K^G(\rho_1-\rho_2))=\delta_S^G(\rho_1-\rho_2)\in Q_S=\mathcal{M}_S$. From the proof of the previous case, we have that $\rho_1|_S-\rho_2|_S$ is a trivial representation of $S$. Since the map $R(K; \mathbb{Q})\to R(S; \mathbb{Q})$ is injective, $\rho_1|_K-\rho_2|_K$ is a trivial representation of $K$. This shows $Q_K\subseteq \mathcal{M}_K$. Together with the easier inclusion $\mathcal{M}_K\subseteq Q_K$, we complete the proof of the proposition. %Let $\displaystyle \sum_i \mu_i\delta_K^G(\rho_i)\in Q_K$, where $\mu_i\in R(K; \mathbb{Q})$ and $\rho_i$ is a representation of $G$. Then $\displaystyle f_2^*\left(\sum_i\mu_i\delta_K^G(\rho_i)\right)=\sum_i\mu_i|_S\delta_S^G(\rho_i)$. By commutativity of the diagram immediately above, $\displaystyle \sum_i\mu_i|_S\delta_S^G(\rho_i)\in Q_S=\mathcal{M}_S$. From the proof of the previous case, we have that $\rho_i|_S$ is a trivial representation of $S$. Since $R(K; \mathbb{Q})\to R(S; \mathbb{Q})$ is injective, $i^*(\rho_i-\text{dim}\rho_i)=0\in R(S; \mathbb{Q})$. Hence $\displaystyle\sum_i\mu_i\delta_K^G(\rho_i)=\sum_i\mu_i\delta_K^G(\rho_i-\text{dim}\rho_i)$ is in $\mathcal{M}_K$ and $Q_K\subseteq \mathcal{M}_K$. Together with the easier inclusion $\mathcal{M_K}\subseteq Q_K$, we prove the proposition.
	%\textcolor{red}{The above proof has gaps; in particular it is not correct in assuming that if $\displaystyle \sum_i\mu_i\delta_S^G(\rho_i)\in Q_S$ then $\delta_S^G(\rho_i)\in Q_S$. Thus the arguments for showing that $Q_S\subseteq\mathcal{M}_S$ needs fixing. One possible way is to study the image of $r_1^*: K_S^*(G/S; \mathbb{Q})\to K_S^*(Z_G(S)_0/S; \mathbb{Q})$, which can be shown to be injective on $\displaystyle\bigwedge\nolimits_{R(S; \mathbb{Q})}\mathcal{L}_S$ if $(G, S)$ is isotropy formal. In this way we have that $\mathcal{L}_S$ is a free $R(S; \mathbb{Q})$-module. If further we can show that $(r_1^*)^{-1}(r_1^*(\mathcal{L}_S))=\mathcal{L}_S$, then we have $f(\mathcal{L}_S)$ is a $\text{rank }G-\text{rank }S$-dimensional vector space using the fact that $\text{ker}(f)=I(S; \mathbb{Q})\cdot K_S^*(G/S; \mathbb{Q})$ if $(G, S)$ is isotropy formal. Note that the assertion $(r_1^*)^{-1}(r_1^*(\mathcal{L}_S))=\mathcal{L}_S$ means that the `odd' part of the $K$-theory ring $K_S^*(G/S; \mathbb{Q})$ is generated by $\mathcal{L}_S$ and nothing else. We may possibly show it using $T$-duality: $G/S$ can be regarded as the torus bundle $G\times_TT/S$ over the flag variety $G/T$. To figure out the odd part, we can instead consider the twisted $K$-theory $K^*(G/T\times T/S, \tau)$ for a suitable twist $\tau$ so that it is isomorphic, through $T$-duality, to $K^*(G/S; \mathbb{Q})$.}
\end{proof}
\begin{corollary}\label{pequaldel}
		If $(G, K)$ is isotropy formal, then $P=\delta^{G/K}(\text{ker}(i^*))$.
\end{corollary}
\begin{proof}
	Consider the following commutative diagram
	\begin{eqnarray}
		\xymatrix{K_K^*(G/K; \mathbb{Q})\ar[r]^{j_K^*}\ar[d]^f&K_K^*(G; \mathbb{Q})\ar[d]^f\\ K^*(G/K; \mathbb{Q})\ar[r]^{j^*}& K^*(G; \mathbb{Q})}
	\end{eqnarray}
	where both the vertical maps are forgetful maps. By Theorem \ref{RKEF}, the left vertical map is surjective. Together with the commutativity of the diagram, we have $f(Q_K)=Q$. By Lemma \ref{eqSam}, $f(Q_K)=f(\mathcal{M}_K)$, which is $\delta(\text{ker}(i^*))$. As $j^*$ maps $P$ onto $Q$ isomorphically by Proposition \ref{kthrystructure}, we have $P=\delta^{G/K}(\text{ker}(i^*))$ as desired.
	%Let us first prove the case where $K$ is a torus subgroup $S$. Let $N_G(S)$ be the normalizer subgroup of $S$ in $G$ and $Z_G(S)$ the centralizer of $S$ in $G$. Then $(G/S)^S=N_G(S)/S\cong N_G(S)/Z_G(S)\times Z_G(S)/S$. Consider the following commutative diagram
	%\begin{eqnarray}
	%	\xymatrix@+4pc{K_S^*(G/S)\ar[r]^{j_S^*}\ar[d]^{r_1^*}&K_S^*(G)\ar[d]^{r_2^*}\\ K_S^*(N_G(S)/Z_G(S)\times Z_G(S)/S)\ar[r]^{k^*}& K_S^*(Z_G(S))}
	%\end{eqnarray}
	%where $r_1^*$ and $r_2^*$ are restriction to the fixed point sets, and $k^*$ is induced by the map $k: Z_G(S)\to N_G(S)/Z_G(S)\times Z_G(S)/S$ which factors through $Z_G(S)/S$. Upon localization of the various equivariant $K$-theory rings in the diagram as $R(S)$-algebra to algebras over the fraction field of $R(S)$, $r_1^*$ and $r_2^*$ become algebra isomorphism by Segal's Localization Theorem (cf. \cite[Proposition 4.1]{Se2}). Note that the $R(S; \mathbb{Q})$-submodule $\mathcal{N}_S$ of $K_S^*(G; \mathbb{Q})$ generated by the primitive elements $\{\delta_S(\rho)|\rho\in R(G; \mathbb{Q})\}$ is isomorphic to $\Omega^1_{R(G; \mathbb{Q})/\mathbb{Q}}\otimes_{R(G; \mathbb{Q})}R(S; \mathbb{Q})$ by Theorem \ref{kthrygrp}(\ref{BZ}). {Its rank as an $R(S; \mathbb{Q})$-module is $\text{rank }G$} by Proposition \ref{rkkahler}. Let $r_{Z_G(S)_0}^*$ be the restriction map $K_S^*(Z_G(S))\to K_S^*(Z_G(S)_0)$ induced by the inclusion of the identity component $Z_G(S)_0$ of $Z_G(S)$. Then the composition $r_{Z_G(S)_0}^*\circ r_2^*$ sends $\mathcal{N}_S$ into $\mathcal{P}$, the $R(S; \mathbb{Q})$-submodule of the primitive elements of $K_S^*(Z_G(S)_0; \mathbb{Q})$. After localization of $R(S; \mathbb{Q})$ to its fraction field, $r_{Z_G(S)_0}^*\circ r_2^*$ is surjective as $r_{Z_G(S)_0}^*$ is surjective and $r_2^*$ becomes an isomorphism. Since $K_S^*(Z_G(S)_0; \mathbb{Q})\cong R(S; \mathbb{Q})\otimes K^*(Z_G(S)_0; \mathbb{Q})$, $\mathcal{P}$ is isomorphic to $R(S; \mathbb{Q})\otimes \text{Im}(\delta: R(Z_G(S)_0; \mathbb{Q})\to K^{-1}(Z_G(S)_0; \mathbb{Q}))$. Thus 
	%\begin{align*}
	%	\text{rank}_{R(S; \mathbb{Q})}\mathcal{P}&=\text{dim}_\mathbb{Q}\text{Im}(\delta: R(Z_G(S)_0; \mathbb{Q})\to K^{-1}(Z_G(S)_0; \mathbb{Q}))\\
	%									&=\text{rank }Z_G(S)_0\\
	%									&=\text{rank }G
	%\end{align*}
	%It follows that $r_{Z_G(S)_0}^*\circ r_2^*$ restricted to $\mathcal{N}_S$ is an isomorphism after localization. 
\end{proof}
\begin{remark}
The converse of Corollary \ref{pequaldel} is not true. Consider the pair $(SU(6), SU(3)\times SU(3))$ where $SU(3)\times SU(3)$ is embedded into $SU(6)$ as block diagonal matrices. This pair is known to be not formal (cf. \cite[pp. 486--488]{GHV}). So by Proposition \ref{kthrystructure}, $\text{dim }P<\text{rank }SU(6)-\text{rank }SU(3)\times SU(3)=1$ and hence $P=0$. By Lemma \ref{key}, $\delta^{G/K}(\text{ker}(i^*))=0$ as well. However, the pair is not isotropy formal (cf. \cite[Example 3.5]{CFok}).
\end{remark}
\begin{lemma}
	The $R(K; \mathbb{Q})$-module $\mathcal{L}_K$ is of rank $\text{rank}G-\text{rank}K$.
\end{lemma}
\begin{proof}
	%Since both the horizontal maps and the right vertical maps are injective, so is the left vertical map $j_S^*$ (when restricted to $\mathcal{L}_S$). It follows that $\mathcal{L}_S$ is a free $R(S; \mathbb{Q})$-module because it is isomorphic to its image which is a $R(S; \mathbb{Q})$-submodule of the free module $K_S^*(G; \mathbb{Q})$. Now 
	Consider the conormal bundle sequence 
		\begin{eqnarray}\label{conormseq}\text{ker}(i^*)/(\text{ker}(i^*))^2\longrightarrow \Omega^1_{R(G; \mathbb{Q})}\otimes_{R(G; \mathbb{Q})}R\longrightarrow \Omega^1_{R/\mathbb{Q}}\longrightarrow 0.\end{eqnarray}
	Tensoring with $R(S; \mathbb{Q})$ over $R$, we have
	\begin{eqnarray*}\text{ker}(i^*)/\text{ker}(i^*)^2\otimes_R R(K; \mathbb{Q})\longrightarrow\Omega^1_{R(G; \mathbb{Q})/\mathbb{Q}}\otimes_{R(G; \mathbb{Q})}R(K; \mathbb{Q})\longrightarrow\Omega^1_{R/\mathbb{Q}}\otimes_R R(K; \mathbb{Q})\longrightarrow 0.\end{eqnarray*}
	The image of the first map, which is the $R(K; \mathbb{Q})$-module generated by $\{\delta_K^G(\rho)|\rho\in\text{ker}(i^*)\}$, is isomorphic to $\mathcal{L}_K$. %which sends $(\rho+\text{ker}(i^*)^2)\otimes 1$ to $d\rho\otimes 1$, can be identified with $j_S^*|_{\mathcal{L}_S}$. So the above sequence can be rewritten as the following short exact sequence 
	%\[0\longrightarrow \mathcal{L}_S\longrightarrow\Omega^1_{R(G; \mathbb{Q})/\mathbb{Q}}\otimes_{R(G; \mathbb{Q})}R(S; \mathbb{Q})\longrightarrow \Omega^1_{R/\mathbb{Q}}\otimes_R R(S; \mathbb{Q})\longrightarrow 0.\]
	Localizing further to the fraction field of $R(K; \mathbb{Q})$ and invoking exactness, we have
	\[\text{rank}\mathcal{L}_K=\text{rank}_{R(K; \mathbb{Q})}\Omega^1_{R(G; \mathbb{Q})/\mathbb{Q}}\otimes_{R(G; \mathbb{Q})}R(K; \mathbb{Q})-\text{rank}_{R(K; \mathbb{Q})}\Omega^1_{R/\mathbb{Q}}\otimes_RR(K; \mathbb{Q}).\]
	Here $\text{rank}\Omega^1_{R(G; \mathbb{Q})/\mathbb{Q}}\otimes_{R(G; \mathbb{Q})}R(K; \mathbb{Q})=\text{rank}G$ by Proposition \ref{rkkahler}, and $\text{rank}\Omega^1_{R/\mathbb{Q}}\otimes_RR(K; \mathbb{Q})=\text{rank}_R\Omega^1_{R/\mathbb{Q}}=\text{Krull dim} R=\text{Krull dim} R(K; \mathbb{Q})=\text{rank}K$, with the first equality due to $R$ being a subring of $R(K; \mathbb{Q})$, the second one due to \cite[Lemma 7.13]{CFok} and the last one due to the going-up theorem and $R(K; \mathbb{Q})$ being integral over $R$ (cf. \cite[Remark, p.60]{AM} and \cite[Proposition 3.2]{Se}). All in all, $\mathcal{L}_K$ is of rank $\text{rank}G-\text{rank}K$
\end{proof}
%\begin{lemma}
%	Let $S$ be a compact torus and act on a finite CW-complex $X$ equivariantly formally. Let $r_S^*: K_S^*(X; \mathbb{Q})\to K_S^*(X^S; \mathbb{Q})\cong R(S; \mathbb{Q})\otimes K^*(X^S; \mathbb{Q})$ be the restriction map. Then for any nonzero $K$-theory class $a\in K^*(X^S; \mathbb{Q})$, $f((r_S^*)^{-1}(R(S; \mathbb{Q})\otimes a))$ is nonzero in $K^*(X; \mathbb{Q})$.
%\end{lemma}

%\begin{lemma}\label{kvolumeform}
%	Let $\pi: G\to\text{pt}$ be the collapsing map. There exists a nonzero element $\displaystyle a\in\bigwedge\nolimits_{\mathbb{Q}}^{\text{rank}G}(\text{Im}(\delta)\otimes\mathbb{Q})$ such that $\pi_*(a)=1$. 
%\end{lemma}
%\begin{proof}
%	There exists a compact Lie group $\widetilde{G}$ which is a finite cover of $G$ such that it is the product of a torus $T$ and a simply-connected, connected Lie group $G_0$. Let $\widetilde{\pi}: \widetilde{G}\to\text{pt}$ be the collapsing map for $\widetilde{G}$, $\{e_1, e_2, \cdots, e_{\text{dim}T}\}$ a basis of the character lattice of $T$, $\rho_1, \cdots, \rho_{\text{rank}G_0}$ the fundamental representations of $G_0$. Then by \cite[Proposition 1 and Equation (3)]{A}, we have $\displaystyle\widetilde{\pi}_*\left(\prod_{i=1}^{\text{dim}T}\delta(e_i)\prod_{j=1}^{\text{rank}G_0}\delta(\rho_j)\right)=1$. Let $p: \widetilde{G}\to G$ be the projection map, and define $\displaystyle a:=p_*\left(\prod_{i=1}^{\text{dim}T}\delta(e_i)\prod_{j=1}^{\text{rank}G_0}\delta(\rho_j)\right)$. Then 
%	\begin{align*}
%		\pi_*(a)&=\pi_*\left(p_*\left(\prod_{i=1}^{\text{dim}T}\delta(e_i)\prod_{j=1}^{\text{rank}G_0}\delta(\rho_j)\right)\right)\\
%			&=\widetilde{\pi}_*\left(\prod_{i=1}^{\text{dim}T}\delta(e_i)\prod_{j=1}^{\text{rank}G_0}\delta(\rho_j)\right)\\
%			&=1.
%	\end{align*}
%	Moreover, 
%	\begin{align*}
%		p^*(a)&=p^*\circ p_*\left(\prod_{i=1}^{\text{dim}T}\delta(e_i)\prod_{j=1}^{\text{rank}G_0}\delta(\rho_j)\right)\\
%			&=|\widetilde{G}/G|\prod_{i=1}^{\text{dim}T}\delta(e_i)\prod_{j=1}^{\text{rank}G_0}\delta(\rho_j).
%	\end{align*}
%	Note that $p^*: K^*(G; \mathbb{Q})\to K^*(\widetilde{G}; \mathbb{Q})$ is an isomorphism of exterior algebras. Thus 
%	\[a=(p^*)^{-1}\left(|\widetilde{G}/G|\prod_{i=1}^{\text{dim}T}\delta(e_i)\prod_{j=1}^{\text{rank}G_0}\delta(\rho_j)\right)\in\bigwedge\nolimits^{\text{dim}T+\text{rank}G_0}\text{Im}(\delta)\otimes\mathbb{Q}=\bigwedge\nolimits^{\text{rank}G}\text{Im}(\delta)\otimes\mathbb{Q}.\]
%\end{proof}

%\begin{proof}[Proof of Theorem \ref{formality}]
	%Again we first consider the case where $K$ is a subtorus $S$. Assume that the pair $(G, S)$ is isotropy formal. Let $N_G(S)$ be the normalizer subgroup of $S$ in $G$. Note that $N_G(S)/Z_G(S)$ is a subgroup of the Weyl group of $G$ with respect to a maximal torus containing $S$, and $(G/S)^S=N_G(S)/S\cong N_G(S)/Z_G(S)\times Z_G(S)/S\cong\coprod_{w\in N_G(S)/Z_G(S)}wZ_G(S)/S$, where $Z_G(S)$ is connected. Consider 
%\begin{eqnarray}
%	\xymatrix{K_S^*((G/S)^S; \mathbb{Q})\ar[r]^{r_*}&K_S^*(G/S; \mathbb{Q})\ar[r]^{r^*}& K_S^*((G/S)^S; \mathbb{Q}).}
%\end{eqnarray}

%Here $r^*$ is the restriction to fixed points, and $r_*$ is the $K$-theoretic pushforward induced by in the inclusion of fixed points. The latter makes sense because the normal bundle $N$ of $(G/S)^S$ and $G/S$ is isomorphic to the complex vector bundle $\displaystyle \coprod_{w\in N_G(S)/Z_G(S)}wZ_G(S)/S\times\bigoplus_{\substack{\alpha\in R_+\\ \alpha|_{\mathfrak{s}}\neq 0}}\mathbb{C}_{w\alpha|_{\mathfrak{s}}}$ (here $R_+$ is the set of positive roots on a Cartan subalgebra $\mathfrak{t}$ containing the Lie algebra $\mathfrak{s}$ of $S$), making the inclusion map equivariantly $K$-orientable.  We also have that $G/S$ is $\text{Spin}^c$ $G$-equivariantly (and hence $S$-equivariantly). That is because the $G$-equivariant third integral Stiefel-Whitney class of $G/S$ is 0 as it lives in $H_G^3(G/S; \mathbb{Z})\cong H_S^3(\text{pt}, \mathbb{Z})$, which is 0. The fixed point set $(G/S)^S$ is $S$-equivariantly $\text{Spin}^c$ as well since it is a disjoint union of Lie groups, whose tangent bundles are trivial. Let $\pi_{G/S*}$ and $\pi_{(G/S)^S*}$ be the pushforward induced by collapsing $G/S$ and $(G/S)^S$ respectively. If we $K$-orient $r_*$ by the $\text{Spin}^c$ structure given by the complex structure of $N^*$, then we have $\pi_{G/S*}\circ r_*=\pi_{(G/S)^S*}$. By Lemma \ref{kvolumeform}, we can find an element $\displaystyle a\in\bigwedge\nolimits^{\text{rank}Z_G(S)/S}_{\mathbb{Q}}(\text{Im}(\delta)\otimes\mathbb{Q})\subseteq K^*(Z_G(S)/S; \mathbb{Q})$ such that $\pi_{Z_G(S)/S*}(a)$ is nonzero. Let $\underline{a}$ be the element in 
%\[K_S^*((G/S)^S; \mathbb{Q})\cong\bigoplus_{w\in N_G(S)/Z_G(S)}R(S; \mathbb{Q})\otimes K^*(wZ_G(S)/S; \mathbb{Q})\cong\bigoplus_{w\in N_G(S)/Z_G(S)} R(S; \mathbb{Q})\otimes K^*(Z_G(S)/S; \mathbb{Q})\]
%(the last isomorphism is due to the obvious diffeomorphism between $wZ_G(S)/S$ and $Z_G(S)/S$) such that its $w$-component $\underline{a}_w$ is $1\otimes a$ for all $w\in N_G(S)/Z_G(S)$. Let $b=r_*(\underline{a})$. Then 
%\begin{align*}
%	\pi_{G/S*}(b)&=\pi_{G/S*}\circ r_*(\underline{a})\\
%			&=\pi_{(G/S)^S*}(\underline{a})\\
%			&=\sum_{i=1}^{|N_G(S)/Z_G(S)|}\pi_{Z_G(S)/S*}(a)\\
%			&=|N_G(S)/Z_G(S)|.
%\end{align*}
%Moreover, 
%\begin{align*}
%	(r^*(b))_w&=(e_S(N^*)\cdot\underline{a})_w\\
%		&=e_S(N^*)_w\cdot \underline{a}_w\\
%		&=\prod_{\substack{\alpha\in R_+\\ \alpha|_{\mathfrak{s}}\neq 0}}(1-t^{-w\alpha|_{\mathfrak{s}}})\otimes a.
%\end{align*}
%We partition the set $R_+$ into two subsets, namely $\mathcal{A}:=\{\alpha\in R_+|\ \alpha|_{\mathfrak{s}}\neq 0, w\cdot\alpha\in R_+\text{ for all }w\in N_G(S)/Z_G(S)\}$ and $\mathcal{B}:=R_+\setminus\mathcal{A}$. Let $s^\alpha$ denote the $K$-theory class of the homogeneous vector bundle $G\times_S\mathbb{C}_\alpha$. Since the latter becomes the vector bundle $Z_G(S)\times_S\mathbb{C}_\alpha\cong Z_G(S)/S\times\mathbb{C}_\alpha$ upon restriction to $Z_G(S)/S$, we have $r^*(s^\alpha)_w=t^{w\alpha}\otimes 1$. Let $D$ be the $R(S; \mathbb{Q})$-submodule of $K_S^*((G/S)^S; \mathbb{Q})$ consisting of elements whose $w$-components are the same for all $w\in N_G(S)/Z_G(S)$ and lie in $R(S; \mathbb{Q})\otimes(\delta(R(Z_G(S)/S))\otimes\mathbb{Q})$. Define $c:=\frac{b}{\prod_{\alpha\in\mathcal{B}}(1-s^{-\alpha|_{\mathfrak{s}}})}$, which is in the localization of $K_S^*(G/S; \mathbb{Q})$. On the one hand, $r^*(c)\in D$ because
%\begin{align*}
%	r^*(c)_w&=\frac{r^*(b)_w}{r^*\left(\prod_{\alpha\in\mathcal{B}}(1-s^{-\alpha|_{\mathfrak{s}}})\right)_w}\\
%		&=\frac{\prod_{\substack{\alpha\in R_+\\ \alpha|_{\mathfrak{s}}\neq 0}}(1-t^{-w\alpha|_{\mathfrak{s}}})\otimes a}{\prod_{\alpha\in\mathcal{B}}(1-t^{-w\alpha|_{\mathfrak{s}}})}\\
%		&=\prod_{\alpha\in\mathcal{A}}(1-t^{-w\alpha|_{\mathfrak{s}}})\otimes a\\
%		&=\prod_{\alpha\in\mathcal{A}}(1-t^{-\alpha|_{\mathfrak{s}}})\otimes a.
%\end{align*}
%One the other hand, note that $r^*(\mathcal{L}_S)\subseteq D$. Moreover, $D\cong R(S; \mathbb{Q})\otimes(\delta(R(Z_G(S)/S))\otimes\mathbb{Q})$ and 
%\begin{align*}
%	\text{rank}_{R(S; \mathbb{Q})}D&=\text{rank}Z_G(S)/S\ (\text{By Proposition }\ref{repringreg})\\
%					&=\text{rank}Z_G(S)-\text{rank}S\\
%					&=\text{rank}G-\text{rank}S\\
%					&=\text{rank}_{R(G; \mathbb{Q})}\mathcal{L}_S.
%\end{align*}
%By the Segal's localization theorem (\cite[Proposition 4.1]{Se2}), $r^*$ maps $(\mathcal{L}_S)_{(0)}$ isomorphically to $R(S; \mathbb{Q})_{(0)}\otimes(\delta(R(Z_G(S)/S))\otimes\mathbb{Q})$, and hence $\displaystyle\bigwedge\nolimits_{R(S; \mathbb{Q})_{(0)}}^{\text{rank}G-\text{rank}S}\mathcal{L}_S$ is isomorphic to $R(S; \mathbb{Q})_{(0)}\otimes a$ through $r^*$.  Thus $\displaystyle c\in\bigwedge\nolimits_{R(S; \mathbb{Q})_{(0)}}^{\text{rank}G-\text{rank}S}\mathcal{L}_S$. It follows, from Corollary \ref{pequaldel} and the fact from Proposition \ref{kthrystructure} that the ordinary $K$-theroy classes of homogeneous vector bundles lie in $R(S; \mathbb{Q})/i^*I(G; \mathbb{Q})$, that $f(b)\in R(S; \mathbb{Q})/i^*I(G; \mathbb{Q})\otimes\bigwedge_{\mathbb{Q}}^{\text{rank}G-\text{rank}S}P$. As 
%\[\pi_{G/S*}(f(b))=f(\pi_{G/S*}(b))=f(|N_G(S)/Z_G(S)|)=|N_G(S)/Z_G(S)|\neq 0, \]
%$f(b)$ is also nonzero. Thus in particular, $\bigwedge^{\text{rank}G-\text{rank}S}P\neq 0$ and $\text{dim}P\geq \text{rank}G-\text{rank}S$. By Proposition \ref{kthrystructure}, $\text{dim}P=\text{rank}G-\text{rank}S$ and we have $G/S$ is formal as desired.
%By the equivariant formality of the isotropy action on $G/S$ and \cite[Theorem A.4]{RK}, the map $r^*$
	%and 
	%\[r_G^*: K_S^*(G; \mathbb{Q})\to K_S^*(Z_G(S); \mathbb{Q})\cong R(S; \mathbb{Q})\otimes K^*(Z_G(S); \mathbb{Q})\] 
	%is injective.  %and thus both $K_S^*(G/S; \mathbb{Q})$ and $K_S^*(G; \mathbb{Q})$ are free $R(S; \mathbb{Q})$-modules. 
	%We also have that the map $R(Z_G(S)/S; \mathbb{Q})\to R(Z_G(S); \mathbb{Q})$ is injective and the image is generated by the (actual) representations of $Z_G(S)$ which restrict to trivial representations of $S$. Thus $\ell_S^*$ maps $\delta(R(Z_G(S)/S))\subset K^*(Z_G(S)/S; \mathbb{Q})$ injectively to $\delta(R(Z_G(S)))\subset K^*(Z_G(S); \mathbb{Q})$. %Note that $r^*_{G/S}$ sends $\mathcal{L}_S$ to $R(S; \mathbb{Q})\otimes 1\otimes\delta(R(Z_G(S)/S))$ and thus $\mathcal{L}_S\subseteq (r^*_{G/S})^{-1}(R(S; \mathbb{Q})\otimes 1\otimes\delta(R(Z_G(S)/S)))$.
	
	% The Segal localization theorem (\cite[Proposition 4.1]{Se2}) implies that the rank of $(r_{G/S}^*)^{-1}(R(S; \mathbb{Q})\otimes 1\otimes\delta(R(Z_G(S)/S)))$ as a $R(S; \mathbb{Q})$-module is the rank of $R(S; \mathbb{Q})\otimes 1\otimes\delta(R(Z_G(S)/S))$, which is $\text{dim}_\mathbb{Q}\delta(R(Z_G(S)/S))$, which in turn is $\text{rank}Z_G(S)/S$ by Remark \ref{deltaredux} and Proposition \ref{repringreg}. The latter, to be more precise, is $\text{rank}Z_G(S)-\text{rank}S=\text{rank}G-\text{rank}S$. For simplicity let us denote $(r_{G/S}^*)^{-1}(R(S; \mathbb{Q})\otimes 1\otimes\delta(R(Z_G(S)/S)))$ by $\mathcal{Z}_S$. By restricting the top row of the above diagram to the primitive parts of the equivariant $K$-theory, we have the following commutative diagram. 
	%\begin{eqnarray}
	%	\xymatrix{\mathcal{Z}_S\ar[r]^{r_{G/S}^*}\ar[d]^{j_S^*}& R(S; \mathbb{Q})\otimes 1\otimes \delta(R(Z_G(S)/S))\ar[d]^{\ell^*_S}\\ K_S^*(G; \mathbb{Q})\ar[r]^{r_G^*}& K_S^*(Z_G(S); \mathbb{Q})}
	%\end{eqnarray}
	%Since the two horizontal maps and the right vertical map are injective, so is the the left vertical map $j_S^*$ (restricted to $\mathcal{Z}_S$). Note that $\ell_S^*(R(S; \mathbb{Q})\otimes 1\otimes\delta(R(Z_G(S)/S)))$ is a module summand of the free module $K_S^*(Z_G(S); \mathbb{Q})$. Let $\mathcal{Y}_S$ be a submodule of $K_S^*(Z_G(S); \mathbb{Q})$ which is complement to $\ell_S^*(R(S; \mathbb{Q})\otimes 1\otimes\delta(R(Z_G(S)/S)))$. By the injectivity of the various maps in the previous commutative diagram, we have
	%\[K_S^*(G; \mathbb{Q})\cong j_S^*(\mathcal{Z}_S)\oplus (r_G^*)^{-1}(\mathcal{Y}_S).\]
	%Consider another commutative diagram as follows.
	%\begin{eqnarray}
	%	\xymatrix{\mathcal{Z}_S\ar[r]^{j_S^*}\ar[d]_f&K_S^*(G; \mathbb{Q})\ar[d]^f\\ K^*(G/S; \mathbb{Q})\ar[r]^{j^*}& K^*(G; \mathbb{Q})}
	%\end{eqnarray}
	%\begin{claim}
	%	The $R(S; \mathbb{Q})$-module $\mathcal{L}_S$ is free of rank $\text{rank}G-\text{rank}S$.
	%\end{claim}
	%Since both the horizontal maps and the right vertical maps are injective, so is the left vertical map $j_S^*$ (when restricted to $\mathcal{L}_S$). It follows that $\mathcal{L}_S$ is a free $R(S; \mathbb{Q})$-module because it is isomorphic to its image which is a $R(S; \mathbb{Q})$-submodule of the free module $K_S^*(G; \mathbb{Q})$. Now consider the conormal bundle sequence 
	%	\begin{eqnarray}\label{conormseq}\text{ker}(i^*)/(\text{ker}(i^*))^2\longrightarrow \Omega^1_{R(G; \mathbb{Q})}\otimes_{R(G; \mathbb{Q})}R\longrightarrow \Omega^1_{R/\mathbb{Q}}\longrightarrow 0.\end{eqnarray}
	%Tensoring with $R(S; \mathbb{Q})$ over $R$, we have
	%\begin{eqnarray*}\text{ker}(i^*)/\text{ker}(i^*)^2\otimes_R R(S; \mathbb{Q})\longrightarrow\Omega^1_{R(G; \mathbb{Q})/\mathbb{Q}}\otimes_{R(G; \mathbb{Q})}R(S; \mathbb{Q})\longrightarrow\Omega^1_{R/\mathbb{Q}}\otimes_R R(S; \mathbb{Q})\longrightarrow 0.\end{eqnarray*}
	%The first term of the sequence in fact is isomorphic to $\mathcal{L}_S$ through the map $\rho\mapsto \delta_S^{G/S}(\rho)$ for $\rho\in\text{ker}(i^*)$, and the first map, which sends $(\rho+\text{ker}(i^*)^2)\otimes 1$ to $d\rho\otimes 1$, can be identified with $j_S^*|_{\mathcal{L}_S}$. So the above sequence can be rewritten as the following short exact sequence 
	%\[0\longrightarrow \mathcal{L}_S\longrightarrow\Omega^1_{R(G; \mathbb{Q})/\mathbb{Q}}\otimes_{R(G; \mathbb{Q})}R(S; \mathbb{Q})\longrightarrow \Omega^1_{R/\mathbb{Q}}\otimes_R R(S; \mathbb{Q})\longrightarrow 0.\]
	%Localizing further to the fraction field of $R(S; \mathbb{Q})$ and invoking exactness, we have
	%\[\text{rank}\mathcal{L}_S=\text{rank}_{R(S; \mathbb{Q})}\Omega^1_{R(G; \mathbb{Q})/\mathbb{Q}}\otimes_{R(G; \mathbb{Q})}R(S; \mathbb{Q})-\text{rank}_{R(S; \mathbb{Q})}\Omega^1_{R/\mathbb{Q}}\otimes_RR(S; \mathbb{Q}).\]
	%Here $\text{rank}\Omega^1_{R(G; \mathbb{Q})/\mathbb{Q}}\otimes_{R(G; \mathbb{Q})}R(S; \mathbb{Q})=\text{rank}G$ by Proposition \ref{rkkahler}, and $\text{rank}\Omega^1_{R/\mathbb{Q}}\otimes_RR(S; \mathbb{Q})=\text{rank}_R\Omega^1_{R/\mathbb{Q}}=\text{Krull dim} R=\text{Krull dim} R(S; \mathbb{Q})=\text{rank}S$, with the first equality due to $R$ being a subring of $R(S; \mathbb{Q})$, the second one due to \cite[Lemma 7.13]{CFok} and the last one due to the going-up theorem and $R(S; \mathbb{Q})$ being integral over $R$ (cf. \cite[Remark, p.60]{AM} and \cite[Proposition 3.2]{Se}). All in all, $\mathcal{L}_S$ and the isomorphic $R(S; \mathbb{Q})$-module $\text{ker}(i^*)/(\text{ker}(i^*))^2\otimes_R R(S; \mathbb{Q})$ are free $R(S; \mathbb{Q})$-modules of rank $\text{rank}G-\text{rank}S$. 
	
	%By \cite[Lemma 3.3]{CFok}, $\text{ker}(i^*)/(\text{ker}(i^*))^2$ is a free $R$-module of the same rank. Any elements in $\text{ker}(i^*)/(\text{ker}(i^*))^2$, through the isomorphism $\text{ker}(i^*)/(\text{ker}(i^*))^2\otimes_RR(S; \mathbb{Q})\cong\mathcal{L}_S$, are of the form $\delta_S^{G/S}(\rho)$ for $\rho\in\text{ker}(i^*)$. Let $\{\delta_S^{G/S}(\rho_i)|\rho_i\in\text{ker}(i^*), 1\leq i\leq \text{rank}G-\text{rank}S\}$ be a $R$-module basis of $\text{ker}(i^*)/(\text{ker}(i^*))^2$. So it is also a $R(S; \mathbb{Q})$-module basis of $\mathcal{L}_S$. %Consider the following commutative diagram.
	%\begin{eqnarray}
	%	\xymatrix@+6.0pc{\mathcal{L}_S\ar[r]^{r_{G/S}^*}\ar[d]_f& R(S; \mathbb{Q})\otimes 1\otimes\delta(\mathcal{K}_S)\subseteq K_S^*((G/S)^S; \mathbb{Q})\ar[d]^f\\ P\ar[r]^{r^*}& \delta(\mathcal{K}_S)\subseteq K^*((G/S)^S; \mathbb{Q})}
	%\end{eqnarray}
	%As $r_{G/S}^*(\delta_S^{G/S}(\rho_i))=1\otimes 1\otimes \delta(\rho_i|_{Z_G(S)})\in K_S^*((G/S)^S; \mathbb{Q})$ and $r_{G/S}^*$ is injective, $\{1\otimes 1\otimes \delta(\rho_i|_{Z_G(S)})|\rho_i\in\text{ker}(i^*), 1\leq i\leq \text{rank}G-\text{rank}S\}$ is linearly independent over $R(S; \mathbb{Q})$ in $R(S; \mathbb{Q})\otimes 1\otimes\delta(\mathcal{K}_S)$. In fact, it is a module basis of $R(S; \mathbb{Q})\otimes 1\otimes \delta(\mathcal{K}_S)$ as the rank of the former is $\text{dim}_\mathbb{Q}\delta(\mathcal{K}_S)=\text{dim}_\mathbb{Q}\delta(R(Z_G(S)/S))=\text{rank}Z_G(S)-\text{rank}S=\text{rank}G-\text{rank}S$. Thus $r_{G/S}^*|_{\mathcal{L}_S}$ is an isomorphism. The left vertical map is onto by Corollary \ref{pequaldel}, and so is the right vertical map. By commutativity, we have $r^*(P)=\delta(\mathcal{K}_S)$, meaning that $\text{dim}_\mathbb{Q}P\geq\text{dim}_\mathbb{Q}\delta(\mathcal{K}_S)=\text{rank}G-\text{rank}S$. But $P$ is spanned by $\{\delta^{G/S}(\rho_i)|\rho_i\in\text{ker}(i^*), 1\leq i\leq \text{rank}G-\text{rank}S\}$. So $\text{dim}_\mathbb{Q}P=\text{rank}G-\text{rank}S$ and by Proposition \ref{kthrystructure}, $G/S$ is formal as desired. 
	%Applying the forgetful map $f$, we get the set $\{\delta^{G/S}(\rho_i)| \rho_i\in\text{ker}(i^*), 1\leq i\leq\text{rank}G-\text{rank}S\}$ and claim that it is a $\mathbb{Q}$-basis of $f(\mathcal{L}_S)$, which is $P$ by Corollary \ref{pequaldel}. What remains to be shown is that the set is linearly independent. Suppose on the contrary it is linearly dependent, i.e. there exist $r_i\in\mathbb{Q}$, not all zero, such that $\sum_i r_i\delta^{G/S}(\rho_i)=0$. 
	
	%Now we apply $\cdot\otimes_{R(S; \mathbb{Q})}\mathbb{Q}$ to the injective map $j_S^*: \mathcal{L}_S\to\Omega^1_{R(G; \mathbb{Q})/\mathbb{Q}}\otimes_{R(G; \mathbb{Q})}R(S; \mathbb{Q})$ and get the map $j_S^*\otimes_{R(S; \mathbb{Q})}\text{Id}_\mathbb{Q}: \mathcal{L}_S\otimes_{R(S; \mathbb{Q})}\mathbb{Q}\to \Omega^1_{R(G; \mathbb{Q})}\otimes_{R(G; \mathbb{Q})}\mathbb{Q}$. Equivariant formality of the $S$-action on $G/S$ implies, by Theorem \ref{RKEF}, that the forgetful map induces an isomorphism $K_S^*(G/S; \mathbb{Q})\otimes_{R(S; \mathbb{Q})}\mathbb{Q}\to K^*(G/S; \mathbb{Q})$. Thus $\mathcal{L}_S\otimes_{R(S; \mathbb{Q})}\mathbb{Q}\cong f(\mathcal{L}_S)=\delta^{G/S}(\text{ker}(i^*))$, which is $P$ by Corollary \ref{pequaldel}. Thus the map $j_S^*\otimes_{R(S; \mathbb{Q})}\text{Id}_\mathbb{Q}$ can be identified with the map $j^*: P\to\Omega^1_{R(G; \mathbb{Q})/\mathbb{Q}}\otimes_{R(G; \mathbb{Q})}\mathbb{Q}$, which is injective by Proposition \ref{kthrystructure}. It follows that $\text{dim}_\mathbb{Q}P=\text{rank}\mathcal{L}_S=\text{rank}G-\text{rank}S$, which amounts to $G/S$ being formal by Theorem \ref{kthrystructure}.
	
%Since the image under $\ell_S^*$ of $R(S; \mathbb{Q})\otimes 1\otimes\delta(\mathcal{K}_S)$ is a summand submodule of $R(S; \mathbb{Q})\otimes \delta(R(Z_G(S)))$, and 
%	\[\mathcal{L}_S=(r_{G/S}^*)^{-1}(R(S; \mathbb{Q})\otimes 1\otimes\delta(\mathcal{K}_S))\cap (j_S^*)^{-1}(\Omega^1_{R(G; \mathbb{Q})/\mathbb{Q}}\otimes_{R(G; \mathbb{Q})}R(S; \mathbb{Q})), \]
%	(the image under $j_S^*$ of) $\mathcal{L}_S$ is a summand submodule of the $R(S; \mathbb{Q})$-module $\Omega^1_{R(G; \mathbb{Q})/\mathbb{Q}}\otimes_{R(G; \mathbb{Q})}R(S; \mathbb{Q})$ which is free. Let $\{a_1, \cdots, a_{\text{rank}G-\text{rank}S}\}$ be a module basis of $\mathcal{L}_S$. Then it can be extended to a module basis $\{a_1, \cdots, a_{\text{rank}G-\text{rank}S}, \cdots, a_{\text{rank}G}\}$ of $\Omega^1_{R(G; \mathbb{Q})/\mathbb{Q}}\otimes_{R(G; \mathbb{Q})}R(S; \mathbb{Q})$ (note that by Proposition \ref{rkkahler} the rank of $\Omega^1_{R(G; \mathbb{Q})/\mathbb{Q}}\otimes_{R(G; \mathbb{Q})}R(S; \mathbb{Q})$ is $\text{rank}G$). Applying the forgetful map $f$ to $\Omega^1_{R(G; \mathbb{Q})/\mathbb{Q}}\otimes_{R(G; \mathbb{Q})}R(S; \mathbb{Q})$ gives $\Omega^1_{R(G; \mathbb{Q})/\mathbb{Q}}\otimes_{R(G; \mathbb{Q})}\mathbb{Q}$ whose dimension is $\text{rank}G$ by Corollary \ref{repringreg}. Thus $\{f(a_1), \cdots, f(a_{\text{rank}G-\text{rank}S}), \cdots, f(a_{\text{rank}G})\}$ is linearly independent, and so is $\{f(a_1), \cdots, f(a_{\text{rank}G-\text{rank}S})\}$. The latter spans $f(\mathcal{L}_S)$, which by Corollary \ref{pequaldel} is $P$. Thus $\text{dim}P=\text{rank}G-\text{rank}S$ and by Proposition \ref{kthrystructure}, $G/S$ is formal.
		%Thus $\mathcal{L}_S\subseteq (r_{G/S}^*)^{-1}(R(S; \mathbb{Q})\otimes 1\otimes \delta(\mathcal{K}_S))$. Note that both $\mathcal{L}_S$ and $(r_{G/S}^*)^{-1}(R(S; \mathbb{Q})\otimes 1\otimes \delta(\mathcal{K}_S))$ are free $R(S; \mathbb{Q})$-modules of rank $\text{rank}G-\text{rank}S$. Fix module bases for both free modules and let $A\in R(S; \mathbb{Q})$ be the determinant of the change of bases matrix.%Now we would like to show that $(r^*)^{-1}(R(S; \mathbb{Q})\otimes 1\otimes\delta(\mathcal{K}_S))=\mathcal{L}_S$. %Note that $r^*$ becomes an algebra isomorphism upon localizing the $K$-theory rings as $R(S; \mathbb{Q})$-algebras at the zero ideal of $R(S; \mathbb{Q})$ by Segal's Localization Theorem (cf. \cite[Proposition 4.1]{Se2}), and 
	%\begin{align*}
	%	\text{rank}_{R(S; \mathbb{Q})}R(S; \mathbb{Q})\otimes 1\otimes\delta(R(Z_G(S)_0/S; \mathbb{Q}))&=\text{dim}_\mathbb{Q}\delta(R(Z_G(S)_0/S; \mathbb{Q}))\\
	%	&=\text{dim}_\mathbb{Q}\Omega^1_{R(Z_G(S)_0/S)/\mathbb{Q}}\otimes_{R(Z_G(S)_0/S; \mathbb{Q})}\mathbb{Q}\\
	%	&=\text{rank}Z_G(S)_0/S\ (\text{by Proposition }\ref{repringreg})\\
	%	&=\text{rank}G-\text{rank}S\\
	%	&=\text{rank}_{R(S; \mathbb{Q})}\mathcal{L}_S.
	%\end{align*}
	%It follows that $r^*$ restricts to an isomorphism between $(\mathcal{L}_S)_{(0)}$ and $R(S; \mathbb{Q})_{(0)}\otimes 1\otimes\delta(R(Z_G(S)_0/S; \mathbb{Q}))$. 
	
	%Suppose for the sake of contradiction that $G/S$ is not formal. Then $\text{dim}P<\text{rank}G-\text{rank }S$. By Corollary \ref{pequaldel}, the forgetful map $f: K_S^*(G/S; \mathbb{Q})\to K^*(G/S; \mathbb{Q})$ maps $\mathcal{L}_S$ onto $P$ and thus $\bigwedge\nolimits_{R(S; \mathbb{Q})}^*\mathcal{L}_S$ onto $\bigwedge\nolimits_\mathbb{Q}^*P$. The top degree part $\bigwedge_{R(S; \mathbb{Q})}^{\text{rank}G-\text{rank}S}\mathcal{L}_S$ is nonzero as $\text{rank}_{R(S; \mathbb{Q})}\mathcal{L}_S=\text{rank}G-\text{rank}S$, and lies in $\text{ker}(f)$ because $\bigwedge\nolimits^{\text{rank}G-\text{rank}S}_\mathbb{Q}P=0$. By \cite[Lemma 2.4]{F2} and isotropy formality of $(G, S)$, $\bigwedge\nolimits_{R(S; \mathbb{Q})}^{\text{rank}G-\text{rank}S}\mathcal{L}_S\subseteq I(S; \mathbb{Q})\cdot K_S^*(G/S; \mathbb{Q})$. The map $r^*$ sends $\bigwedge\nolimits_{R(S; \mathbb{Q})}^{\text{rank}G-\text{rank}S}(\mathcal{L}_S)_{(0)}$ isomorphically to $\bigwedge\nolimits_{R(S; \mathbb{Q})}^{\text{rank}G-\text{rank}S}R(S; \mathbb{Q})_{(0)}\otimes 1\otimes\delta(R(Z_G(S)_0/S; \mathbb{Q}))=R(S; \mathbb{Q})_{(0)}\otimes 1\otimes\bigwedge\nolimits^{\text{rank}G-\text{rank}S}_\mathbb{Q}\delta(R(Z_G(S)_0/S; \mathbb{Q}))$, which should be in $r^*(I(S; \mathbb{Q})\cdot K_S^*(G/S; \mathbb{Q})_{(0)})=I(S; \mathbb{Q})_{(0)}\otimes K^*(N_G(S)/Z_G(S)\otimes\pi_0(Z_G(S)); \mathbb{Q})\otimes K^*(Z_G(S)_0/S; \mathbb{Q})$. That is a contradiction, however. Thus $\text{dim}_\mathbb{Q}P=\text{rank}G-\text{rank}S$ and by Proposition \ref{kthrystructure}, $G/S$ is formal.
	
%	The general case $G/K$ with $(G, K)$ isotropy formal can be reduced to the previous case $G/S$ where $S$ is a maximal torus of $K$. That is because
%	\begin{enumerate}
%		\item $(G, K)$ is isotropy formal if and only if $(G, S)$ is (cf. \cite[Theorem 1.1]{Ca}), 
%		\item the kernels of both restriction maps $R(G; \mathbb{Q})\to R(S; \mathbb{Q})$ and $R(G; \mathbb{Q})\to R(K; \mathbb{Q})$ are the same, as the former map factors through $R(K; \mathbb{Q})$ and the restriction $R(K; \mathbb{Q})\to R(S; \mathbb{Q})$ is injective, and 
%		\item for $G/K$ (resp. $G/S$), $P=\delta^{G/K}(\text{ker}(i^*))$ (resp. $P=\delta^{G/S}(\text{ker}(i^*))$ (Corollary \ref{pequaldel}).
%	\end{enumerate}
%	Thus for $G/K$, $\text{dim}_\mathbb{Q}P=\text{dim}_\mathbb{Q}\delta^{G/K}(\text{ker}(i^*))=\text{dim}_\mathbb{Q}\delta^{G/S}(\text{ker}(i^*))=\text{rank}G-\text{rank}S=\text{rank}G-\text{rank}K$, and it is formal by Proposition \ref{kthrystructure}.
%\end{proof}	
%\begin{remark}
%From the above proof, we see that whether $(G, S)$ is isotropy formal or not, we always have the exact sequence
%\[\mathcal{L}_S\longrightarrow\Omega^1_{R(G; \mathbb{Q})/\mathbb{Q}}\otimes_{R(G; \mathbb{Q})}R(S; \mathbb{Q})\longrightarrow \Omega^1_{R/\mathbb{Q}}\otimes_R R(S; \mathbb{Q})\longrightarrow 0.\]
%By exactness and the analysis of the ranks of the second and third terms of the above sequence in the previous proof, we have $\text{rank}_{R(S; \mathbb{Q})}\mathcal{L}_S\geq\text{rank}G-\text{rank}S$. On the other hand, we have the map $r^*_{G/S}|_{\mathcal{L}_S}:\mathcal{L}_S\to R(S; \mathbb{Q})\otimes 1\otimes\delta(R(Z_G(S)/S))\subseteq K_S^*((G/S)^S; \mathbb{Q})$ and that $r_{G/S}^*$ is an isomorphism after localizing at the zero ideal of $R(S; \mathbb{Q})$ by the Segal localization theorem (\cite[Proposition 4.1]{Se2}). Thus 
%\[\text{rank}_{R(S; \mathbb{Q})}\mathcal{L}_S\leq \text{rank}_{R(S; \mathbb{Q})}R(S; \mathbb{Q})\otimes 1\otimes\delta(R(Z_G(S)/S))=\text{dim}_\mathbb{Q}\delta(R(Z_G(S)/S))\]
%with the last term being $\text{rank}G-\text{rank}S$ from the previous proof. All in all, $\mathcal{L}_S$ has rank $\text{rank}G-\text{rank}S$ in general and $\bigwedge\nolimits^{\text{rank}G-\text{rank}S}_{R(S; \mathbb{Q})}\mathcal{L}_S\neq 0$.
%\end{remark}

%\begin{theorem}
%	Assume the same conditions as in Theorem \ref{rep} and let $S$ be a torus subgroup of $G$. The pair $(G, S)$ is formal if and only if there exists $a\in\bigwedge^{\text{rank }G-\text{rank }S}P$ such that $\pi_*(a)\neq 0$, where $\pi: G/S\to\text{pt}$ is the collapsing map.
%\end{theorem}
%\begin{proof}
%	If there exists $\displaystyle a\in\bigwedge\nolimits^{\text{rank }G-\text{rank }S}P$ such that $\pi_*(a)\neq 0$, then $a\neq 0$ and thus $\text{dim }P\geq\text{rank }G-\text{rank }S$. By Proposition \ref{kthrystructure}, $\text{dim }P$ is $\text{rank }G-\text{rank }S$ and thus $(G, S)$ is formal. On the other hand, if $(G, S)$ is formal, then $\text{dim }P=\text{rank }G-\text{rank }S$ by Proposition \ref{kthrystructure}. We shall show that any nonzero element $a\in\bigwedge\nolimits^{\text{rank }G-\text{rank }S}P$ satisfies $\pi_*(a)\neq 0$.
%\end{proof}
	\begin{proposition}\label{rankSam}
		The pair $(G, K)$ is isotropy formal if and only if $\text{dim }\delta^{G/K}(\text{ker}(i^*))=\text{rank }G-\text{rank }K$.
	\end{proposition}
	\begin{proof}Suppose $(G, K)$ is isotropy formal. Then by \cite[Theorem A]{CFok}, $G/K$ is formal and by Proposition \ref{kthrystructure}, $\text{dim }P=\text{rank }G-\text{rank }K$. By Corollary \ref{pequaldel}, $\text{dim }\delta^{G/K}(\text{ker}(i^*))=\text{rank }G-\text{rank }K$. Conversely, suppose that $\text{dim }\delta^{G/K}(\text{ker}(i^*))=\text{rank }G-\text{rank }K$. Then by Proposition \ref{kthrystructure} and Lemma \ref{key}, 
	\[\text{rank }G-\text{rank }K=\text{dim }\delta^{G/K}(\text{ker}(i^*))\leq \text{dim }P\leq \text{rank }G-\text{rank }K.\]
	Thus $P=\delta^{G/K}(\text{ker}(i^*))$ and $\text{dim }P=\text{rank }G-\text{rank }K$. By Proposition \ref{kthrystructure} again, we have the isomorphism $K^*(G/K; \mathbb{Q})\cong \left(R(K; \mathbb{Q})/i^*I(G; \mathbb{Q})\right)\otimes\bigwedge\nolimits^*P$. The first tensor factor consists of $K$-theory classes represented by associated vector bundles $G\times_K V$, which admits the equivariant structure via the left $K$-action on $G$. The second tensor factor admits equivariant lifts as well by the fact that the $K$-theory forgetful map $K^*_K(G/K; \mathbb{Q})\to K^*(G/K; \mathbb{Q})$ maps $\delta^{G/K}_K(\text{ker}(i^*))$ onto $P$. We have that the forgetful map is onto and $(G, K)$ is indeed isotropy formal by Theorem \ref{RKEF}. 
	\end{proof}
	%In fact one can use the proof of the above Proposition to strengthen Corollary \ref{pequaldel} to the following
	%\begin{proposition}
	%	The pair $(G, K)$ is isotropy formal if and only if $P=\delta^{G/K}(\text{ker}(i^*))$.
	%\end{proposition}
	\begin{proof}[Proof of Theorem \ref{rep}]
	By Proposition \ref{rankSam}, it remains to show that $\text{dim }\delta^{G/K}(\text{ker}(i^*))=\text{rank }G-\text{rank }K$ if and only if $R:=i^*R(G; \mathbb{Q})$ is regular at the ideal $I:=i^*I(G; \mathbb{Q})$. By \cite[Corollary 3.7]{ABC+}, $R$ is regular at $I$ if and only if 
	\[\text{dim}_\mathbb{Q}(R/I)\otimes_R \Omega^1_{R/\mathbb{Q}}=\text{Krull dim }R_I, \text{where }R/I\cong\mathbb{Q}.\]
	Since $R$ is an integral domain finitely generated over $\mathbb{Q}$ and $I$ is a maximal ideal, the Krull dimension of $R_I$ equals the Krull dimension of $R$ (cf. \cite[Corollary 11.27]{AM}), which in turn is also the Krull dimension of $R(K; \mathbb{Q})$ since $R(K; \mathbb{Q})$ is integral over $R$ (cf. \cite[Proposition 3.2]{Se}, \cite[Remark, p.60]{AM}). By \cite[Lemma 7.12]{CFok}, the Krull dimension of $R(K)$ is rank $K$. Now we have shown another
	\begin{claim}
		The ring $R$ is regular at the ideal $I$ if and only if $\text{dim}_\mathbb{Q}(R/I)\otimes_R\Omega^1_{R/\mathbb{Q}}=\text{rank }K$. 
	\end{claim}
	To finish the proof, we shall show that 
	\begin{claim}
		$\text{dim}_\mathbb{Q}(R/I)\otimes_R\Omega^1_{R/\mathbb{Q}}=\text{rank }K$ if and only if $\text{dim }\delta(\text{ker}(i^*))=\text{rank }G-\text{rank }K$.
	\end{claim} 
	Note that in the conormal bundle sequence (\ref{conormseq}), the first map sends $\rho+(\text{ker}(i^*))^2$ to $d\rho\otimes 1$. Tensoring the sequence with $\cdot\otimes_R R/I$ yields the following sequence of $\mathbb{Q}$-vector spaces.
	\[(\text{ker}(i^*)/(\text{ker}(i^*)^2))\otimes_R R/I\longrightarrow\Omega^1_{R(G; \mathbb{Q})/\mathbb{Q}}\otimes_{R(G; \mathbb{Q})}\mathbb{Q}\longrightarrow\Omega^1_{R/\mathbb{Q}}\otimes_R R/I\longrightarrow 0.\]
	The image of the first map of the last sequence is isomorphic to $d(\text{ker}(i^*))\otimes_{R(G; \mathbb{Q})}\mathbb{Q}$, which in turn is isomorphic to $\delta(\text{ker}(i^*))$ by Remark \ref{deltaredux} and to $\delta^{G/K}(\text{ker}(i^*))$ by the isomorphism $j^*$ (Proposition \ref{kthrystructure}). By exactness we have 
	\[\text{dim}_\mathbb{Q}\delta^{G/K}(\text{ker}(i^*))+\text{dim}_\mathbb{Q}\Omega^1_{R/\mathbb{Q}}\otimes_R R/I=\text{dim}_\mathbb{Q}\Omega^1_{R(G; \mathbb{Q})/\mathbb{Q}}\otimes_{R(G; \mathbb{Q})}\mathbb{Q}.\]
	By Proposition \ref{repringreg}, the RHS of the above equation is $\text{rank }G$. 
	%It remains to show that 
	%\[\displaystyle\text{dim}_\mathbb{Q} \Omega^1_{R(G; \mathbb{Q})/\mathbb{Q}}\otimes_{R(G; \mathbb{Q})}\mathbb{Q}=\text{rank }G,\] 
	Now the proof of Theorem \ref{rep} is complete. 
\end{proof}
\begin{corollary}\label{regrep}
	The representation ring $R(G; \mathbb{Q})$ is regular at the augmentation ideal $I(G; \mathbb{Q})$. 
\end{corollary}
\begin{proof}
	We know that $(G, G)$ is always isotropy formal as the $G$-action on $G/G$ is trivial. By Theorem \ref{rep}, $R(G; \mathbb{Q})$ is regular at its augmentation ideal, and so is $R(G)$. 
\end{proof}
\begin{remark}
	In fact, the following string of equalities
	\begin{align*}
		&\text{dim}_\mathbb{Q}\Omega^1_{R(G; \mathbb{Q})/\mathbb{Q}}\otimes_{R(G; \mathbb{Q})}\mathbb{Q}\\
		=&\text{rank}G\ (\text{by Proposition \ref{repringreg}})\\
		=&\text{Krull dim} R(G; \mathbb{Q})\ (\text{by \cite[Lemma 7.12]{CFok}})\\
		=&\text{Krull dim} R(G; \mathbb{Q})_{I(G; \mathbb{Q})}\ (\text{by \cite[Corollary 11.27]{AM}})
	\end{align*} 	
	also imply that $R(G; \mathbb{Q})$ is regular at $I(G; \mathbb{Q})$ by \cite[Corollary 3.7]{ABC+}
\end{remark}
\begin{example}
	We will show by explicit computation that $R(\text{PSU}(3))$ is regular at the augmentation ideal, in accordance with our result and thus correct a claim made in \cite{Ho}. The representation ring $R(\text{PSU}(3))$ is isomorphic to 
	\[\mathbb{Z}[x, y, z]/(y^3-y^2-xz-2y(x+z)-x-y-z), \]
	where $x=[V_{3L_1}]$, $y=[V_{2L_1+L_2}]$ and $z=[V_{3L_1+3L_2}]$. As a check, one may make the linear change of variables $x= X+1-2Y$, $y=Y-1$ and $z=Z+1-2Y$ and obtain the more compact presentation in \cite[Lemma 7.1]{BZ}:
	\[R(\text{PSU}(3))\cong \mathbb{Z}[X, Y, Z]/(Y^3-XZ).\]
	Note that $\text{dim }V_{3L_1}=\text{dim }V_{3L_1+3L_2}=10$ and $\text{dim }V_{2L_1+L_2}=8$. If we let $\overline{x}=x-10$, $\overline{y}=y-8$ and $\overline{z}=z-10$, then 
	\[R(\text{PSU}(3))\cong\mathbb{Z}[\overline{x}, \overline{y}, \overline{z}]/(\overline{y}^3+23\overline{y}^2-2(\overline{x}+\overline{z})\overline{y}-\overline{x}\overline{z}-27(\overline{x}+\overline{z})+135\overline{y}).\]
	It follows that 
	\[I(\text{PSU}(3))/I(\text{PSU}(3))^2=\text{span}_\mathbb{Z}\{\overline{x}, \overline{y}, \overline{z}\}/\text{span}_\mathbb{Z}\{-27(\overline{x}+\overline{z})+135\overline{y}\}\]
	whose rank is 2, which is equal to the the Krull dimension of $R(\text{PSU}(3))$ (or equivalently the rank of $\text{PSU}(3)$ by \cite[Lemma 7.12(ii)]{CFok}). Thus $R(\text{PSU}(3))$ is indeed regular at its augmentation ideal. Moreover $I(\text{PSU}(3))/I(\text{PSU}(3))^2$ is an (non-free) abelian group of rank 2 instead of a free abelian group of rank 3 as claimed after \cite[Proof of Lemma 4.2]{Ho}. 
\end{example}

%\begin{lemma}
%	Let $\displaystyle p_x(a)=c_na^n+c_{n-1}a^{n-1}+\cdots+c_0-xa^i$ and $q_y(a)=d_ma^m+d_{m-1}a^{m-1}+\cdots+d_0-ya^j$, where $0\leq i\leq n$, $0\leq j\leq m$. Suppose $p_0(a)$ and $q_0(a)$ have $k$ common roots (counting multiplicities). Then the resultant of $p_x(a)$ and $q_y(a)$ is a polynomial in $x$ and $y$ with the lowest degree being $k$. 
%\end{lemma}

We are now ready to demonstrate the first application of Theorem \ref{rep} in this paper, which is to show Theorem \ref{circle}, a characterization of the isotropy formality of $(G, S)$ where $S$ is a circle subgroup, in terms of the self-duality of representations of $G$ when restricted to $S$.
\begin{proof}[Proof of Theorem \ref{circle}] Let $\rho_1, \rho_2, \cdots, \rho_m$ be representations which generate $R(G)$, and $\displaystyle i^*\rho_k=\sum_{n}a_{k, n}t^n\in R(S)$, where $t$ is one-dimensional standard representation of $S$. Suppose that any representation of $G$, when restricted to $S$, is self-dual. Then $a_{k, n}=a_{k, -n}$ for all $k$ and $n$. If we denote $t+t^{-1}-2$ by $s$, then $i^*\rho_k$ can be rewritten as a polynomial $p_k(s)$ of $s$. The regularity of $R:=i^*R(G; \mathbb{Q})$ at the augmentation ideal $i^*I(G; \mathbb{Q})=R\cap (t-1)$ amounts to the smoothness of $\text{Spec} R$ at the point $i^*I(G; \mathbb{Q})$. This is also equivalent to the smoothness of the image of the map $i_\mathbb{R}: \text{Spec} R(S; \mathbb{R})=\text{Spec}\mathbb{R}[t, t^{-1}]\to \text{Spec} R(G; \mathbb{R})=\text{Spec}\mathbb{R}[x_1, x_2, \cdots, x_{m}]/\mathcal{I}\subseteq\text{Spec}\mathbb{R}[x_1, x_2, \cdots, x_m]$ at the point $i_\mathbb{R}((t-1))=I(G; \mathbb{R})$ where $\mathcal{I}$ is the kernel of the map $\mathbb{R}[x_1, x_2, \cdots, x_m]\to R(G; \mathbb{R})$ which sends $x_i$ to $\rho_i$. In other words, the curve with the parametrization $s\mapsto (p_1(s), p_2(s), \cdots, p_m(s))\in\mathbb{R}^m$, which is a reparametrization of the curve $\displaystyle t\mapsto \left(\sum_n a_{1, n}t^n, \cdots, \sum_n a_{m, n}t^n\right)$, is smooth at the point $s=0$ (note that $t=1$ if and only if $s=0$). Thus it suffices to show that $(p_1'(0), p_2'(0), \cdots, p_m'(0))\neq 0$. Differentiating the equation $\displaystyle p_k(s)=\sum_n a_{k, n}t^n$ with respect to $t$ and applying chain rule to the LHS, we get $p_k'(0)$ as follows.
\begin{align*}
	\frac{d}{ds}p_k(s)\cdot\frac{ds}{dt}&=\sum_n na_{k, n}t^{n-1}\\
	\frac{d}{ds}p_k(s)\cdot (1-t^{-2})&=\sum_{n>0}na_{k, n}(t^{n-1}-t^{-n-1})\\
	\frac{d}{ds}p_k(s)&=\sum_{n>0}na_{k, n}(t^{n-1}+t^{n-3}+\cdots+t^{1-n})\\
	p_k'(0)&=\sum_{n>0}n^2a_{k, n}
\end{align*}
The RHS of the last line is not zero because $a_{k, n}\geq 0$ but there exists $n=n_0>0$ such that $a_{k, n_0}>0$. This shows that $R$ is regular at $i^*I(G; \mathbb{Q})$. 

Now suppose that there exists a representation of $G$ whose restriction to $S$ is not self-dual. Let $U$ be a set of representations of $G$ which generate $R(G; \mathbb{R})$, $\overline{U}$ the set of dual representations, and $U\cup \overline{U}=\{\rho_1, \rho_2, \cdots, \rho_m\}$. In this way $\{\rho_1, \rho_2, \cdots, \rho_m\}$ generates $R(G; \mathbb{Q})$ and has the nontrivial involution of taking dual (If $G$ is in addition simply-connected, the set $(\rho_1, \rho_2, \cdots, \rho_m)$ can be taken to be the set of fundamental representations as it generates $R(G; \mathbb{Q})$ and is invariant under the involution of taking dual). Let this involution send $\rho_k$ to $\rho_{\sigma(k)}$, where $\sigma$ is a nontrivial involution on $\{1, 2, \cdots, m\}$. We want to show that $R$ is not regular at $i^*I(G; \mathbb{Q})$. We have that $i^*\rho_k=q_k(t)$ for some Laurent polynomial $q_k$. By the same reasoning as in the previous paragraph, it suffices to show that the image curve $t\mapsto (q_1(t), q_2(t), \cdots, q_m(t))\subseteq \mathbb{R}^m$ for $t>0$ is not smooth at the point $t=1$. Suppose for the sake of contradiction that the curve is smooth at $t=1$. We shall first show that the map $t\mapsto (q_1(t), q_2(t), \cdots, q_m(t))$ is injective for $t$ in an open neighborhood of 1. Choose $k$ such that $k\neq \sigma(k)$. It suffices to show that the map $t\mapsto(q_k(t), q_{\sigma(k)}(t))$ is injective in an open neighborhood of 1. As $q_k(t)$ is a nonconstant Laurent polynomial, there exists $\delta_1>0$ such that $q'_k(t)>0$ or $q'_k(t)<0$ for $t\in(1, 1+\delta_1)$. So $q_k$ and hence $t\mapsto(q_k(t), q_{\sigma(k)}(t))$ is injective for $t\in(1, 1+\delta_1)$. Similarly, there exists $\delta_2>0$ such that $q_k(t)-q_{\sigma(k)}(t)>0$ (resp. $q_k(t)-q_{\sigma(k)}(t)<0$) for $t\in(1, 1+\delta_2)$. Take $\delta=\min\{\delta_1, \delta_2\}$. For $\displaystyle t\in\left(\frac{1}{1+\delta}, 1\right)$, $t^{-1}\in(1, 1+\delta)$, and the map $t\mapsto(q_k(t), q_{\sigma(k)}(t))=(q_{\sigma(k)}(t^{-1}), q_k(t^{-1}))$ is injective and $q_{\sigma(k)}(t^{-1})-q_k(t^{-1})<0$ (resp. $q_{\sigma(k)}(t^{-1})-q_k(t^{-1})>0$). Thus $t\mapsto (q_k(t), q_{\sigma(k)}(t))$ is injective for $\displaystyle t\in\left(\frac{1}{1+\delta}, 1+\delta\right)$. We may further restrict $\displaystyle\left(\frac{1}{1+\delta}, 1+\delta\right)$ to a smaller neighborhood of 1 so that the part of the curve for $t$ in that neighborhood is smooth. Now we can reparametrize that part of the curve by $(r_1(s), r_2(s), \cdots, r_m(s))$ so that the tangent vectors with respect to $s$ is nowhere vanishing (we may parametrize by arc length, for example). In particular the tangent vector at the point $(q_1(1), q_2(1), \cdots, q_m(1))=(\text{dim}\rho_1, \text{dim}\rho_2, \cdots, \text{dim}\rho_m)$ is nonzero. Let $s=s_0$ correspond to the point $(\text{dim}\rho_1, \text{dim}\rho_2, \cdots, \text{dim}\rho_m)$. Note that, since $G$ is semisimple, for any $k$ $q_k(t)$ is of the form $\displaystyle\sum_n a_{k, n}t^n$ where $\displaystyle\sum_n na_{k, n}=0$ and $a_{k, n}\geq 0$. By the AM-GM inequality, 
\[\displaystyle\sum_n a_{k, n}t^n\geq\left(\sum_n a_{k, n}\right)\sqrt{\prod_n (t^n)^{a_{k, n}}}=\text{dim}\rho_k,\] 
and equality holds if and only if $t=1$. Thus each coordinate $r_k(s)$ of the curve achieves its minimum at $s=s_0$ and so $r'(s_0)=(r'_1(s_0), r'_2(s_0), \cdots, r'_m(s_0))=(0, 0, \cdots, 0)$, contradicting the nonvanishing of the tangent vector at $(\text{dim}\rho_1, \text{dim}\rho_2, \cdots, \text{dim}\rho_m)$. It follows that the curve is not smooth at the point $(\text{dim}\rho_1, \text{dim}\rho_2, \cdots, \text{dim}\rho_m)$, which more precisely is a cusp given that each coordinate of the curve is bounded below by the corresponding coordinate of the point.
\end{proof}
%\begin{remark}
%\end{remark}
%\begin{proof}[Proof of Theorem \ref{indexthry}]
	
%	If $(G, S)$ is not isotropy formal, then $\text{dim}f(\mathcal{L}_S)=\text{dim}\delta^{G/S}(\text{ker}(i^*))<\text{rank }G-\text{rank }S$ by Corollary \ref{pequaldel} and Proposition \ref{kthrystructure}. It follows that $\displaystyle f\left(\bigwedge\nolimits_{R(S; \mathbb{Q})}^{\text{rank }G-\text{rank }S}\mathcal{L}_S\right)=\bigwedge\nolimits_\mathbb{Q}^{\text{rank }G-\text{rank }S}f(\mathcal{L}_S)=0$, and for any $\displaystyle a\in\bigwedge\nolimits_{R(S; \mathbb{Q})}^{\text{rank }G-\text{rank }S}\mathcal{L}_S$, $f(\pi_*(a))=\pi_*(f(a))=\pi_*(0)=0$. Then we have $\pi_*(a)\in I(S; \mathbb{Q})$ for any $\displaystyle a\in\bigwedge\nolimits_{R(S; \mathbb{Q})}^{\text{rank }G-\text{rank }S}\mathcal{L}_S$.
%\end{proof}
\section{Characterization of isotropy formality of corank 1 isotropy actions}
\begin{theorem}\label{invtpoly}
	Let $G$ be a compact connected Lie group of rank $n$, $S$ a torus subgroup of dimension being $n-1$, $T$ a maximal torus of $G$ containing $S$, $\mu\in \chi(T)$ a primitive character of $T$ such that $S=\text{ker }\mu$, $L$ a primitive weight in the weight lattice $\mathfrak{t}^*_\mathbb{Z}$ corresponding to $\mu$ after exponentiation, and $W$ the Weyl group of $G$. Define
	\[\Phi_L:=\left(\prod_{w\in W/W_{[L]}} wL\right)^{n(L)}\in \text{Sym}^*(\mathfrak{t}^*)^W,\]
	where $W_{[L]}$ is the stabilizer of $[L]:=\{\pm L\}\in \mathfrak{t}^*/\{\pm 1\}$, $n(L)=1$ or 2 depending on whether $\displaystyle\prod_{w\in W/W_{[L]}}wL$ is $W$-invariant or $W$-anti-invariant, i.e. 
	\[\displaystyle w'(\prod_{w\in W/W_{[L]}}wL)=\text{sgn}(w')\prod_{w\in W/W_{[L]}}wL\text{ for all }w'\in W.\] 
	Then $(G, S)$ is isotropy formal if and only if $\Phi_L\notin (f_1, \cdots, f_n)^2$, where $f_1, \cdots, f_n$ are the fundamental $W$-invariant homogeneous polynomials on $\mathfrak{t}$, i.e. $\text{Sym}^*(\mathfrak{t}^*)^W=\mathbb{R}[f_1, \cdots, f_n]$.
\end{theorem}
\begin{proof}
	First, we use $\widehat{R}$ to denote the completion of $R$ at its augmentation ideal. Recall that the augmentation ideal of $\text{Sym}^*(V)$ is $\text{Sym}^{\geq 1}(V)$. Then with the identification $\text{Sym}^*(\mathfrak{t}^*)^W=\mathbb{R}[f_1, \cdots, f_n]$, the augmentation ideal of $\text{Sym}^*(\mathfrak{t}^*)^W$ is $(f_1, \cdots, f_n)$. Consider the following diagram.
	\begin{eqnarray*}
		\xymatrix{\widehat{R(G; \mathbb{R})}\ar[r]^{\text{ch}}\ar[d]^{i^*}& \widehat{\text{Sym}^*(\mathfrak{t}^*)^W}\ar[d]^{i^*}\\ \widehat{R(S; \mathbb{R})}\ar[r]^{\text{ch}}&\widehat{\text{Sym}^*(\mathfrak{s}^*)}}
	\end{eqnarray*}
	The horizontal maps are Chern character maps from the completed equivariant $K$-theory coefficient ring (identified with the completed representation ring) to the completed equivariant cohomology coefficient ring (identified with the completed polynomial ring on the Lie algebra of the acting group). By \cite[Theorem 5.3]{CFok}, the Chern character map is a ring isomorphism. The diagram commutes due to the functoriality of the Chern character. By Theorem \ref{rep}, $(G, S)$ is isotropy formal if and only if $i^*R(G; \mathbb{R})$ is regular at $i^*I(G; \mathbb{R})$. The latter condition then is equivalent to $i^*\widehat{R(G; \mathbb{R})}$ being regular at $i^*I(G; \mathbb{R})$, which in turn is equivalent to $\text{ker}(i^*: \widehat{R(G; \mathbb{R})}\to\widehat{R(S; \mathbb{R})})\nsubseteq I(G; \mathbb{R})^2$. Applying the Chern character and commutativity of the above diagram, we have that the last condition is equivalent to $\mathfrak{p}:=\text{ker}(i^*: \widehat{\text{Sym}^*(\mathfrak{t}^*)^W}\to \widehat{\text{Sym}^*(\mathfrak{s}^*)})\nsubseteq (f_1, \cdots, f_n)^2$. Note that
	\begin{align*}
		\text{Krull dim }\widehat{\text{Sym}^*(\mathfrak{t}^*)^W}&=\text{Krull dim }\widehat{\text{Sym}^*(\mathfrak{t}^*})\\
													&(\text{because }\widehat{\text{Sym}^*(\mathfrak{t}^*)}\text{ is integral over }\widehat{\text{Sym}^*(\mathfrak{t}^*)^W}\text{ by \cite[Exercise 5.12]{AM}})\\
													&=\text{dim }T\\
													&=\text{rank }G\\
													&=n\\
		\text{Krull dim }\widehat{\text{Sym}^*(\mathfrak{s}^*)}&=\text{dim }S\\
												&=n-1\\
		\text{Krull dim }i^*\widehat{\text{Sym}^*(\mathfrak{t}^*)^W}&=\text{Krull dim }i^*\widehat{R(G; \mathbb{R})}\\
													&=\text{Krull dim }\widehat{R(S; \mathbb{R})}\\
													&(\text{because }\widehat{R(S; \mathbb{R})}\text{ is integral over }i^*\widehat{R(G; \mathbb{R})}\text{ by \cite[Proposition 3.2]{Se}})\\
													&=\text{dim }S\\
													&=n-1.
	\end{align*}
	It follows that $\mathfrak{p}$ is a prime ideal of height 1 of $\widehat{\text{Sym}^*(\mathfrak{t}^*)^W}$. Since $\widehat{\text{Sym}^*(\mathfrak{t}^*)^W}$ is a Noetherian UFD, $\mathfrak{p}$ is a principal ideal (\cite[Lemma 10.120.6]{SP}). Let $\Phi_L$ be a generator of $\mathfrak{p}$. Now isotropy formality of $(G, S)$ is equivalent to $\Phi_L\nsubseteq (f_1, \cdots, f_n)^2$. To finish the proof, it remains to show that $\Phi_L$ can be chosen to be the polynomial in the proposition. Note that 
	\[\mathfrak{p}=(\Phi_L)\subseteq \text{ker}(i^*: \widehat{\text{Sym}^*(\mathfrak{t}^*)}\to \widehat{\text{Sym}^*(\mathfrak{s}^*)})=(L).\]
	The vanishing set $V(L)$ is the hyperplane $\text{ker }L$ in $\mathfrak{t}$. Since $\Phi_L$ is $W$-invariant, $V(\Phi_L)$ should contain $V(L)$ and be $W$-invariant in $\mathfrak{t}$. Thus $V(\Phi_L)$ is the $W$-orbit of the hyperplane $\text{ker }L$, i.e. 
	\[V(\Phi_L)=\bigcup_{w\in W/W_{[L]}}\text{ker }wL.\]
	Thus $\Phi_L$ should be a $W$-invariant polynomial of minimal degree on $\mathfrak{t}$ which vanishes on $\displaystyle \bigcup_{w\in W/W_{[L]}}\text{ker }wL$. The polynomial $\displaystyle \left(\prod_{w\in W/W_{[L]}} wL\right)^{n(L)}$ satisfies these conditions and this completes the proof.
\end{proof}
The following lemma, which is a consequence of Theorem \ref{invtpoly}, stipulates a bound on the size of the Weyl orbit of $[L]$ by the maximum degree of the fundamental $W$-invariant homogeneous polynomials in order for isotropy formality to hold. This places a severe restriction on $L$ which will prove very helpful in the classification of isotropy formality of corank 1 isotropy actions. 
\begin{lemma}\label{degreebound}
	Assume the same conditions as in Theorem \ref{invtpoly}. If $(G, S)$ is isotropy formal, then 
	\[n(L)|W/W_{[L]}|\leq \max\{\text{deg }f_1, \cdots, \text{deg }f_n\}.\]
\end{lemma}
\begin{proof}
	The LHS of the inequality is the (homogeneous) degree of $\Phi_L$. The condition $\Phi_L\nsubseteq (f_1, \cdots, f_n)^2$ is equivalent to $\Phi_L$ having nonzero linear terms of $f_1, \cdots, f_n$ after being rewritten as a polynomial of $f_1, \cdots, f_n$. Thus the degree of $\Phi_L$ must be equal to the degree of one of the fundamental $W$-invariant polynomials. The lemma then follows.
\end{proof}
\begin{proof}[Proof of Theorem \ref{corank1}]
	\begin{enumerate}
		\item $A_n$: In this case, $\displaystyle \mathfrak{t}^*_\mathbb{Z}=\left.\left\{\sum_{i=1}^{n+1}a_i L_i\right| \sum_{i=1}^{n+1}a_i=0\right\}$, $W=S_{n+1}$ which permutes $L_i$, $\text{deg }f_i=i+1$ ($f_i$ is the elementary symmetric polynomials of degree $i+1$) and so the maximum fundamental degree is $n+1$. First, we shall verify that $L=\omega_1=L_1$ gives rise to an isotropy formal action. Note that $\Phi_L=L_1L_2\cdots L_{n+1}$, which is $f_n$. So $\Phi_L\nsubseteq (f_1, \cdots, f_n)^2$ and $(G, S)$ is isotropy formal by Theorem \ref{invtpoly}. When $n=3$, $L=\omega_2=L_1+L_2$ also leads to an isotropy formal action. That is because $\Phi_L=(L_1+L_2)(L_1+L_3)(L_1+L_4)$ (the stabilizer $W_{[L]}$ also include the permutation $(13)(24)$ which sends $L$ to $-L$ and hence stabilizes $[L]$), which is $f_2$ and thus not in $(f_1, \cdots, f_3)^2$.
		
		Now we shall prove that there are no other primitive weights which give rise to isotropy formal actions. Let $\displaystyle L=\sum_{i=1}^{n+1}a_i L_i$ and $m_1, m_2, \cdots, m_r$ the multiplicities of the coefficients $a_1, \cdots, a_{n+1}$. If there are no Weyl group elements which negate $L$, then $\displaystyle |W/W_{[L]}|=\frac{(n+1)!}{m_1!\cdots m_r!}$. If the multiplicity type $\{m_1, \cdots, m_r\}\neq \{1, n\}$, then the multiplicity type with the largest stabilizer is $\{2, n-1\}$. Note that
		\[|W/W_{[L]}|\geq \frac{(n+1)!}{2!(n-1)!}=\binom{n+1}{2}=\frac{(n+1)n}{2}\]
	and the last term is greater than $\max\{\text{deg }f_1, \cdots, \text{deg }f_n\}=n+1$ when $n\geq 3$. So according to Lemma \ref{degreebound}, when $n\geq 3$ and $L$ has the multiplicity type not equal to $\{1, n\}$ and there are no Weyl group elements taking $L$ to $-L$, then $L$ does not give rise to an isotropy formal pair.
		
	If there are Weyl group elements which take $L$ to $-L$, then either the nonzero coefficients of $L$ form pairs of numbers negative of each other, or $L$ is of the form $\sum_{i=1}^{\frac{n+1}{2}}L_i$ when $n+1$ is even. For the former case, the multiplicity type with the largest stabilizer is $\{1, 1, n-1\}$ ($L=L_1-L_2$ is an example with such a multiplicity type). Note that
	\[|W/W_{[L]}|\geq\frac{(n+1)!}{2(n-1)!}=\binom{n+1}{2}=\frac{(n+1)n}{2}\] 
	and again the last term is greater than $\max\{\text{deg }f_1, \cdots, \text{deg }f_n\}=n+1$ when $n\geq 3$. For the latter case, the multiplicity type is $\{\frac{n+1}{2}, \frac{n+1}{2}\}$. Then 
	\[|W/W_{[L]}|=\frac{1}{2}\binom{n+1}{\frac{n+1}{2}}\]
	which is greater than $\max\{\text{deg }f_1, \cdots, \text{deg }f_n\}=n+1$ when $n\geq 5$ and $n$ is odd. 
	
	In sum, using Lemma \ref{degreebound}, we rule out isotropy formality in the following cases.
	\begin{enumerate}
		\item $n\geq 3$, the multiplicity type of $L$ is not $\{1, n\}$ and no Weyl group elements take $L$ to $-L$.
		\item $n\geq 3$, the coefficients of $L$ forms pairs of numbers negative of each other.
		\item $n\geq 5$ and $n$ is odd, and $\displaystyle L=\sum_{i=1}^{\frac{n+1}{2}}L_i$.
	\end{enumerate}
	We are then left with cases $n=1$, 2 and 3. For $n=1$, then $S$ is the trivial subgroup and $(G, S)$ is trivially isotropy formal. For $n=2$ or 3, we can apply the above arguments to show that $L$ has to be one of the weights listed in the table in Theorem \ref{corank1} in order for $(G, S)$ to be isotropy formal.
	\item $B_n$, $C_n$ and $D_n$: The proof of these cases is similar to that of the type $A_n$ case and also involves ruling out the weight $L$ not listed in the table in Theorem \ref{corank1} by showing that the degree of $\Phi_L$ is greater than the maximal degree of fundamental Weyl invariant polynomials, violating the degree bound in Lemma \ref{degreebound}.%$B_n$ and $C_n$ for $n\geq 2$: When $n=2$, $S$ is a circle subgroup, to which Theorem \ref{circle} is applicable. Any representations of compact Lie groups of these two types are self-dual (cf. \cite{Bo}). It follows that their restriction to $S$ are self-dual as well. By Theorem \ref{circle}, $(G, S)$ is isotropy formal in this case. Recall the following Lie theoretic data for $B_n$ and $C_n$, $n\geq 3$.
	%\begin{enumerate}
	%	\item For $B_n$, the weight lattice $\mathfrak{t}^*_\mathbb{Z}$ is $\displaystyle\left.\left\{\sum_{i=1}^n a_i L_i\right|\ (a_1, \cdots, a_n)\in\mathbb{Z}^n\text{ or }\left(\mathbb{Z}+\frac{1}{2}\right)^n\right\}$. The Weyl group is $W=\mathbb{Z}_2^n\rtimes S_n$, where $S_n$ permutes $L_1, \cdots, L_n$ and $(\varepsilon_1, \cdots, \epsilon_n)\in \mathbb{Z}_2^n$ takes $L_i$ to $(-1)^{\varepsilon_i}L_i$. The fundamental degree $\text{deg }f_i$ is $2i$. 
	%\end{enumerate}
	\item $G_2$: In this case the dimension of $S$ is 1. Note that any representations of $G_2$ are self-dual (cf. \cite{Bo}) and so are their restricted representations to $S$. By Theorem \ref{circle}, $(G, S)$ is isotropy formal for any circle subgroup $S$.
	\end{enumerate}
	As an interlude, let us prove a general result about the size of the smallest $W$-orbit of a weight for a general compact, connected simple Lie group $G$, which will serve as the common argument for the non-existence of isotropy formal corank 1 actions for the remaining exceptional cases. Let $\displaystyle L=\sum_{i=1}^n\omega_i$ and $s_i$ be the reflection across the hyperplane orthogonal to the simple root $\alpha_i$. Note that $W$ is generated by $s_i$ for $1\leq i\leq n$. Note that $s_i$ stabilizes $L$ as a weight if and only if $L-\langle L, \alpha_i^\vee\rangle\alpha_i=L$, i.e. $\langle L, \alpha_i^\vee\rangle=0$, or equivalently $a_i=0$. Thus the stabilizer of $L$ as a weight is the subgroup $W_{\{1, \cdots, n\}\setminus\{i|\ a_i=0\}}:=\langle s_i|\ a_i=0\rangle$. Then if the size of the $W$-orbit of $L$ as a weight is the smallest, then only one coefficient among $a_1, \cdots, a_n$ is nonzero. In other words, the $W$-orbits of the fundamental weights are candidates for the smallest orbit. The stabilizer $W_{i}$ of $\omega_i$ is in fact the Weyl group of the maximal parabolic subgroup of the complexified group $G^\mathbb{C}$ corresponding to the Dynkin diagram obtained by deleting the $i$-th node from that of $G$. 
	\begin{enumerate}
		\setcounter{enumi}{3}
		\item $F_4$: The sizes of the $W$-orbits of the fundamental weights of $F_4$ are as follows.
		\begin{align*}
			|W\cdot\omega_1|&=\frac{|W(F_4)|}{|W(B_3)|}=\frac{1152}{2^3\cdot 3!}=24\\
			|W\cdot\omega_2|&=|W\cdot\omega_3|=\frac{|W(F_4)|}{|W(A_1\times A_2)|}=\frac{1152}{2\cdot 6}=96\\
			|W\cdot\omega_4|&=\frac{|W(F_4)|}{|W(C_3)|}=\frac{1152}{2^3\cdot 3!}=24
		\end{align*}
		Thus the $W$-orbits of $\omega_1$ and $\omega_4$ are the smallest. Since $-\text{Id}_{\mathfrak{t}^*}\in W(F_4)$ (cf. \cite{Bo}), $|W(F_4)\cdot[\omega_i]|=\frac{24}{2}=12$ for $i=1, 4$. The degree of $\Phi_{\omega_i}$ for $i=1$ and $4$ is $12\cdot 2=24$ because $\displaystyle\prod_{w\in W/W_{[\omega_i]}}w\omega_i$ is $W$-anti-invariant. However, $\max\{\text{deg }f_1, \cdots, \text{deg }f_4\}=12$. By Lemma \ref{degreebound}, there does not exist any isotropy formal action of corank 1 when $G=F_4$.
		\item $E_6$: The smallest $W$-orbit of a weight of $E_6$ is the orbit of $\omega_1$. Its size is
		\[\frac{|W(E_6)|}{|W(D_5)|}=\frac{51840}{1920}=27.\]
		Since $-\text{Id}_{\mathfrak{t}^*}\notin W(E_6)$, the orbit of $\omega_1$ as a weight is also the $W$-orbit of $[\omega_1]$. The degree of $\Phi_{\omega_1}$ is at least 27 which is greater than $\max\{\text{deg }f_1, \cdots, \text{deg }f_6\}=12$. By Lemma \ref{degreebound}, there does not exist any isotropy formal actions of corank 1 when $G=E_6$.
		\item $E_7$: The smallest $W$-orbit of a weight of $E_7$ is the orbit of $\omega_1$. Its size is
		\[\frac{|W(E_7)|}{|W(E_6)|}=\frac{2903040}{51840}=56.\]
		Since $-\text{Id}_{\mathfrak{t}^*}\in W(E_7)$ (cf. \cite{Bo}), $|W(E_7)\cdot[\omega_1]|=\frac{56}{2}=28$. The degree of $\Phi_{\omega_1}$ is at least 28 which is greater than $\max\{\text{deg }f_1, \cdots, \text{deg }f_7\}=18$. By Lemma \ref{degreebound}, there does not exist any isotropy formal actions of corank 1 when $G=E_7$.
		\item $E_8$: The smallest $W$-orbit of a weight of $E_8$ is the orbit of $\omega_1$. Its size is
		\[\frac{|W(E_8)|}{|W(E_7)|}=\frac{696729600}{2903040}=240.\]
		Since $-\text{Id}_{\mathfrak{t}^*}\in W(E_8)$ (cf. \cite{Bo}), $|W(E_8)\cdot[\omega_1]|=\frac{240}{2}=120$. The degree of $\Phi_{\omega_1}$ is at least 120 which is greater than $\max\{\text{deg }f_1, \cdots, \text{deg }f_8\}=30$. By Lemma \ref{degreebound}, there does not exist any isotropy formal actions of corank 1 when $G=E_8$.
	\end{enumerate}
\end{proof}
\section{Characterization of isotropy formality using index pairing}
Apart from Theorem \ref{rep}, we also characterize isotropy formality in terms of a certain nondegeneracy of the equivariant $K$-theoretic index pairing of two special equivariant $K$-theory classes. This section is devoted to the proof of Theorem \ref{indexthry}, as well as a demonstration of its applicability to two examples.
	\begin{proposition}\label{nondegenformal}
		Let $G$ be a connected compact Lie group and $S$ a torus subgroup of $G$, and $\pi: G/S\to \text{pt}$ the collapsing map. Then $G/S$ is formal if and only if there exist $a\in R(S; \mathbb{Q})/i^*I(G; \mathbb{Q})$ and $b\in \bigwedge\nolimits^{\text{dim }P}P$ such that the $K$-theoretic index pairing $\pi_*(ab)$ is nonzero.
	\end{proposition}
	\begin{proof}
		First, the tangent bundle $TG/S$ is isomorphic to $G\times_S\mathfrak{g}/\mathfrak{s}=G\times_S\left(\mathfrak{t}/\mathfrak{s}\oplus\bigoplus_{\alpha\in R^+}\mathfrak{g}_\alpha^{\mathbb{R}}\right)$, where $\mathfrak{g}_\alpha^{\mathbb{R}}=\left(\mathfrak{g}_\alpha^\mathbb{C}\oplus\mathfrak{g}_{-\alpha}^\mathbb{C}\right)\cap \mathfrak{g}$ (where $\mathfrak{g}_\alpha^\mathbb{C}$ is the root space corresponding to $\alpha$ in $\mathfrak{g}\otimes_\mathbb{R}\mathbb{C}$) for $\alpha\in R^+$ determine an $S$-invariant complex structure of the tangent bundle. Thus $G/S$ is $S$-equivariantly almost complex and hence $S$-equivariantly $\text{Spin}^c$, and it makes sense to consider the $K$-theoretic pushforward $\pi_*$. Suppose $G/K$ is formal. Choose any nonzero element $b$ of $\bigwedge\nolimits^{\text{dim }P}P$. By the nondegeneracy of the index pairing of any compact $\text{Spin}^c$ manifold, there exists $\overline{a}\in K^*(G/S)\otimes\mathbb{Q}$ such that $\pi_*(\overline{a}b)\neq 0$. By Proposition \ref{kthrystructure}, $\overline{a}=a_0+\sum_{i=1}^n a_ib_i$, where $a_i\in R(S; \mathbb{Q})/i^*I(G; \mathbb{Q})$ for $0\leq i\leq n$ and $0\neq b_i\in \bigwedge\nolimits^*P$ for $1\leq i\leq n$. It follows that $\pi_*(a_0b)=\pi_*(\overline{a}b)\neq 0$.
		
		To prove the converse, it suffices to focus on the case where $\text{dim }S<\text{rank }G$ because $G/S$ is formal if $S$ is a maximal torus of $G$. Now we assume that there exists $a\in R(S; \mathbb{Q})/i^*I(G; \mathbb{Q})$ and $b\in \bigwedge\nolimits^{\text{dim }P}P$ such that $\pi_*(ab)\neq 0$. By the Atiyah-Singer index theorem, 
		\[\pi_*(ab)=\int_{G/S}\text{ch}(ab)\text{td}(G/S).\] 
		So $\text{ch}(ab)\text{td}(G/S)$ contains a term which is a ``volume form'' in $H^{\text{dim }G/S}(G/S; \mathbb{Q})$. On the other hand, consider the Cartan algebra $(H^*(BS; \mathbb{Q})\otimes H^*(G; \mathbb{Q}), d)=(\text{Sym}^*(\mathfrak{s}^*)\otimes \bigwedge\nolimits^*V_H, d)$, which is a minimal model for $H^*(G/S; \mathbb{Q})$. Here the differential $d$ satisfies $d(\bigwedge\nolimits^*V_H)\subseteq \text{Sym}^*(\mathfrak{s}^*)$ and $d(\text{Sym}^*(\mathfrak{s}^*))=0$. The exterior degree on $\bigwedge\nolimits^*V_H$ induces an exterior degree on the complex $(\text{Sym}^*(\mathfrak{s}^*)\otimes\bigwedge\nolimits^*V_H)_p:=\text{Sym}^*(\mathfrak{s}^*)\otimes\bigwedge\nolimits^p V_H$ and hence on its cohomology. Let 
		\[P_H:=\{p\in V_H| dp\in \text{Sym}^{\geq 1}(\mathfrak{s}^*)\cdot dV_H\}\]
		and $P'_H$ a graded linear complement of $P_H$ in $V_H$. Then 
		\begin{align*}
			H^*(G/S; \mathbb{Q})&\cong H^*(\text{Sym}^*(\mathfrak{s}^*)\otimes\bigwedge\nolimits^* P'_H, d)\\
							&\cong (H^*(BS; \mathbb{Q})/i^*H^{>0}(BG; \mathbb{Q})\oplus \mathfrak{a})\otimes \bigwedge\nolimits^* P_H
		\end{align*}
		where $\mathfrak{a}$ is the ideal of elements of positive exterior degrees (cf. \cite[pp. 73, 83, 52]{GHV} and \cite[Theorem 2.6]{CFok}). Now suppose for the sake of contradiction that $G/S$ is not formal. Then $\mathfrak{a}\neq 0$ (\cite[Theorem 2.6]{CFok}). Any ``volume form" in $H^*(G/S; \mathbb{Q})$ is of top positive exterior degree and in particular an element of $\mathfrak{a}\otimes \bigwedge\nolimits^{\text{dim }P_H}P_H$. However, the exterior degree of $\text{ch}(ab)\text{td}(G/S)$ is 0. That is because $\text{ch}(b)\in \text{ch}(\bigwedge\nolimits^{\text{dim }P}P)=\bigwedge\nolimits^{\text{dim }P_H}P_H$ and $\text{ch}(a)\in \text{ch}(R(S; \mathbb{Q})/i^*I(G; \mathbb{Q}))=H^*(BS; \mathbb{Q})/i^*H^{>0}(BG; \mathbb{Q})$ are of 0 exterior degree, and so is
		\[\text{td}(G/S)=\prod_{\alpha\in R^+}\frac{c_1(L_\alpha)}{1-e^{-c_1(L_\alpha)}}\in H^*(BS; \mathbb{Q})/i^* H^{>0}(BG; \mathbb{Q})\]
		where $L_\alpha:=G\times_S\mathfrak{g}_\alpha^{\mathbb{R}}=G\times_S\mathbb{C}_\alpha$. This way we have reached a contradiction and completed the proof.
	\end{proof}
	Theorem \ref{indexthry} we are going to prove next can be regarded as the equivariant analogue of Proposition \ref{nondegenformal}.
	\begin{proof}[Proof of Theorem \ref{indexthry}] Suppose $(G, S)$ is isotropy formal. By \cite[Theorem A]{CFok}, $G/S$ is formal. Proposition \ref{nondegenformal} implies that there exist $a\in R(S; \mathbb{Q})/i^* I(G; \mathbb{Q})$ and $b\in \bigwedge\nolimits^{\text{dim }P}P$ such that $\pi_*(ab)\neq 0$. Since $P=\delta^{G/K}(\text{ker}(i^*))$ by Corollary \ref{pequaldel}, there is a natural equivariant lift $\widetilde{b}\in \bigwedge\nolimits^{\text{dim }P}_{R(S; \mathbb{Q})}\mathcal{L}_S$ of $b$. More precisely, $\widetilde{b}\in \bigwedge\nolimits^{\text{rank }G-\text{dim }S}_{R(S; \mathbb{Q})}\mathcal{L}_S$ because $\text{dim }P=\text{rank }G-\text{dim }S$ by Proposition \ref{kthrystructure}. Likewise, there is a natural equivariant lift $\widetilde{a}\in A$ of $a$ because $a$ is represented by the virtual vector bundle constructed using the line bundles $G\times_S\mathbb{C}_\mu$ associated to the representation $\mu$ of $S$, and these line bundles admit the natural $S$-equivariant structure given by the left translation on $G$. By the commutativity of $\pi_*$ and the forgetful map $f$, $f(\pi_*(\widetilde{a}\widetilde{b}))=\pi_*(ab)\neq 0$ and so $\pi_*(\widetilde{a}\widetilde{b})\notin I(S; \mathbb{Q})$ as desired.
	
	Now suppose that there exist $\widetilde{a}\in A$ and $\widetilde{b}\in\bigwedge\nolimits^{\text{rank }G-\text{dim }S}_{R(S; \mathbb{Q})}\mathcal{L}_S$ such that $\pi_*(\widetilde{a}\widetilde{b})\notin I(S; \mathbb{Q})$. Then $\pi_*(f(\widetilde{a})f(\widetilde{b}))\neq 0$. Note that $f(\widetilde{a})\in f(A)=R(S; \mathbb{Q})/i^*I(G; \mathbb{Q})$ and $f(\widetilde{b})\in f(\bigwedge\nolimits^{\text{rank }G-\text{dim }S}_{R(S; \mathbb{Q})}\mathcal{L}_S)=\bigwedge\nolimits^{\text{rank }G-\text{dim }S}_\mathbb{Q}\delta^{G/K}(\text{ker}(i^*))\subseteq \bigwedge\nolimits^{\text{rank }G-\text{dim }S}_\mathbb{Q}P$. Clearly $f(\widetilde{b})\neq 0$, for otherwise $\pi_*(f(\widetilde{a})f(\widetilde{b}))=0$. So $\bigwedge\nolimits^{\text{rank }G-\text{dim }S}P\neq 0$, implying that $\text{dim }P\geq \text{rank }G-\text{dim }S$, but $\text{rank }G-\text{dim }S\leq \text{dim }P$ by Proposition \ref{kthrystructure}. It follows that $\text{dim }P=\text{rank }G-\text{dim }S$ and $\text{dim }\bigwedge\nolimits_\mathbb{Q}^{\text{rank }G-\text{dim }S}P=1$. Thus $\bigwedge\nolimits_\mathbb{Q}^{\text{rank }G-\text{dim }S}\delta^{G/K}(\text{ker}(i^*))=\bigwedge\nolimits_\mathbb{Q}^{\text{rank }G-\text{dim }S}P$ because $\bigwedge\nolimits_\mathbb{Q}^{\text{rank }G-\text{dim }S}\delta^{G/K}(\text{ker}(i^*))$ is a nonzero subspace of $\bigwedge\nolimits^{\text{rank }G-\text{dim }S}_\mathbb{Q}P$ which is one-dimensional. Then we have $P=\delta^{G/K}(\text{ker}(i^*))$ and $\text{dim }\delta^{G/K}(\text{ker}(i^*))=\text{dim }P=\text{rank }G-\text{dim }S$. By Proposition \ref{rankSam}, $(G, S)$ is isotropy formal and this completes the proof.
	\end{proof}
%In this section, we will illustrate our main results (Theorems \ref{rep}, \ref{circle}, \ref{indexthry}) with two examples.
\begin{example}\label{isoex}
	Let $G=SU(3)$, $\displaystyle S=\left.\left\{\begin{pmatrix}t& &\\ &t^{-1}& \\ &&1\end{pmatrix}\right|t\in U(1)\right\}$, $\sigma_{\text{std}}$ the standard representation of $G$. By abuse of notation, we also use $t$ to denote the standard 1-dimensional representation of $S$
	\begin{align*}
		S&\to U(1)\\
		\begin{pmatrix}t& &\\ &t^{-1}& \\ &&1\end{pmatrix}&\mapsto t.
	\end{align*}
	Then $\displaystyle R(G)=\mathbb{Z}[\sigma_{\text{std}}, \bigwedge\nolimits^2\sigma_{\text{std}}]$ and $R(S)=\mathbb{Z}[t, t^{-1}]$. The restriction map $i^*: R(G)\to R(S)$ takes both generators $\sigma_{\text{std}}$ and $\displaystyle\bigwedge\nolimits^2\sigma_{\text{std}}$ to $1+t+t^{-1}$. So $i^*R(G; \mathbb{Q})$ is isomorphic to $\mathbb{Q}[1+t+t^{-1}]$, which is regular at $I=i^*I(G; \mathbb{Q})=(t+t^{-1}-2)$. By Theorem \ref{rep}, $(G, S)$ is an isotropy formal pair. Alternatively, as mentioned before, both restrictions $i^*\sigma_{\text{std}}$ and $i^*\bigwedge\nolimits^2\sigma_{\text{std}}$ are $1+t+t^{-1}$, which is self-dual. By Theorem \ref{circle}, we can also get the isotropy formality of $(G, S)$.
	
	We can also use Theorem \ref{indexthry} to show the isotropy formality of $(G, S)$. Since $\text{ker}(i^*)=(\sigma_{\text{std}}-\bigwedge\nolimits^2\sigma_{\text{std}})$, $\mathcal{L}_S$ is the $R(S; \mathbb{Q})$-submodule of $K^*_S(G/S; \mathbb{Q})$ generated by $b:=\delta^{G/S}_S(\sigma_{\text{std}}(\bigwedge\nolimits^2\sigma_{\text{std}})^{-1})$. We take $a$ to be 1, represented by the equivariant trivial vector bundle over $G/S$. We will show that $\pi_*(ab)=\pi_*(b)=-1$, which is not in $I(S; \mathbb{Q})$. Let $T$ be the maximal torus $\displaystyle\left.\left\{\begin{pmatrix}t_1&&\\ &t_2&\\ && t_3\end{pmatrix}\right| t_i\in U(1), 1\leq i\leq 3, t_1t_2t_3=1\right\}$ of $SU(3)$, which contains $S$. Again by abuse of notation we use $t_i$ to denote the 1-dimensional representation which takes a diagonal matrix from $T$ to the $i$-th diagonal entry. Then fixed point set $(G/S)^S$ of the isotropy action is $T/S\cup gT/S$, where $\displaystyle g=\begin{pmatrix}&-1&\\ 1&&\\ &&1\end{pmatrix}$. The Atiyah-Segal localization formula says that 
	\[\pi_*(b)=\pi_{T/S*}\left(\frac{i_{T/S, G/S}^*(b)}{e_S(N_{T/S}^*)}\right)+\pi_{gT/S*}\left(\frac{i_{gT/S, G/S}^*(b)}{e_S(N^*_{gT/S})}\right).\]
	Here $i_{T/S, G/S}^*$ and $i_{gT/S, G/S}^*$ are restrictions from $K_S^*(G/S)$ to $K_S^*(T/S)$ and $K_S^*(gT/S)$ respectively, $N_{T/S}$ and $N_{gT/S}$ are normal bundles over $T/S$ and $gT/S$ respectively, $\pi_{T/S*}$ and $\pi_{gT/S*}$ are pushforward to $K_S^*(\text{pt})$. The short exact sequence of groups
	\[1{\longrightarrow}S\stackrel{i_{S, T}}{\longrightarrow} T\stackrel{\pi_{T, T/S}}{\longrightarrow}T/S\longrightarrow 1\]
	induces the short exact sequence of character groups
	\[0\longrightarrow \chi(T/S)\stackrel{\pi_{T, T/S}^*}{\longrightarrow} \chi(T)\stackrel{i_{S, T}^*}{\longrightarrow} \chi(S)\longrightarrow 0,\]
	where $i_{S, T}^*t_1=t$, $i_{S, T}^*t_2=t^{-1}$. The maps $i_{S, T}^*$ and $p^*$ extend by linearity to maps $i_{S, T}^*: R(T)\to R(S)$ and $\pi_{T, T/S}^*: R(T/S)\to R(T)$. We have that $R(T/S)=\mathbb{Z}[s, s^{-1}]$ with $\pi_{T, T/S}^*s=t_1t_2$. Consider the commutative diagram
	\[\xymatrix@+4pc{K_S^*(G/S)\ar[r]^{i_{T/S, G/S}^*}\ar[d]_{\pi^*_{G, G/S}}& K_S^*(T/S)\cong R(S)\otimes K^*(T/S)\ar[d]^{\pi^*_{T, T/S}}\\ K_S^*(G)\ar[r]^{i_{T, G}^*}& K_S^*(T)\cong R(S)\otimes K^*(T)}\]
	with maps induced by inclusions and projections. Here $S$ acts on $G$ and $T$ by conjugation. Then
	\begin{align*}
		i_{T}^*\circ\pi^*_{G/S}(b)&=i_T^*(\delta_S^G(\sigma_{\text{std}}-\bigwedge\nolimits^2\sigma_{\text{std}}))\\
							&=i_{S, T}^*t_1\otimes \delta(t_1)+i_{S, T}^*t_2\otimes\delta(t_2)+i_{S, T}^*t_3\otimes\delta(t_3)\\
							&-(i_{S, T}^*t_1t_2\otimes\delta(t_1+t_2)+i_{S, T}^*t_2t_3\otimes\delta(t_2+t_3)+i_{S, T}^*t_3t_1\otimes\delta(t_3+t_1))\\
							&(\text{cf. \cite[Lemma 4.19]{F}})\\
							&=t\otimes\delta(t_1)+t^{-1}\otimes\delta(t_2)+1\otimes\delta(-t_1-t_2)\\
							&-(1\otimes\delta(t_1+t_2)+t^{-1}\otimes\delta(-t_1)+t\otimes\delta(-t_2))\\
							&=(t+t^{-1}-2)\otimes\delta(t_1+t_2).
	\end{align*}
	By the commutativity of the above diagram and the injectivity of $\pi_{T/S}^*$, we have $i_{T/S}^*(b)=(t+t^{-1}-2)\otimes\delta(s)$. Similarly, $i_{gT/S}^*(b)=(t+t^{-1}-2)\otimes\delta(s)\in K_S^*(gT/S)\cong R(S)\otimes K^*(T/S)$. Next, the normal bundle $N_{T/S}$ is a trivial $T$-equivariant complex vector bundle over $T/S$ with weights of the $T$-action on the fibers being the positive roots of the Lie algebra of $G$, i.e. $t_1t_2^{-1}$, $t_2t_3^{-1}$ and $t_1t_3^{-1}$. When restricted to the $S$-action, the weights become $t^2$, $t^{-1}$ and $t$. Thus the $S$-equivariant $K$-theoretic Euler class $e_S(N_{T/S}^*)$ is $(1-t^{-2})(1-t)(1-t^{-1})$. Similarly, we find that $e_S(N_{gT/S}^*)$ is $(1-t^2)(1-t^{-1})(1-t)$. Assembling all these data to the Atiyah-Segal localization theorem yields
	\begin{align*}
		\pi_*(b)&=\pi_{T/S*}\left(\frac{(t+t^{-1}-2)\otimes\delta(s)}{(1-t^{-2})(1-t)(1-t^{-1})}\right)+\pi_{gT/S*}\left(\frac{(t+t^{-1}-2)\otimes\delta(s)}{(1-t^2)(1-t^{-1})(1-t)}\right)\\
				&=\frac{t+t^{-1}-2}{(1-t^{-2})(1-t)(1-t^{-1})}+\frac{t+t^{-1}-2}{(1-t^2)(1-t)(1-t^{-1})}\\
				&=-1
	\end{align*}
	which is not in $I(S; \mathbb{Q})$. By Theorem \ref{indexthry}, we again verify that $(G, S)$ is isotropy formal.
\end{example}

\begin{example}
	Let $G=SU(3)$ and $\displaystyle S=\left.\left\{\begin{pmatrix}t&&\\ &t&\\ &&t^{-2}\end{pmatrix}\right|t\in U(1)\right\}$. The isotropy formality of $(G, S)$ is discussed in \cite[Example 7.19]{CFok}. It is shown there that $i^*R(G; \mathbb{Q})$ is not regular at the ideal $i^*I(G; \mathbb{Q})$ and hence by Theorem \ref{rep}, $(G, S)$ is not isotropy formal. Alternatively, since $i^*\sigma_{\text{std}}=t+t+t^{-2}$ is not self-dual, we have again the non-isotropy formality by Theorem \ref{circle}.
	
	As in Example \ref{isoex}, we will also confirm the non-isotropy formality of $(G, S)$ by using Theorem \ref{indexthry}. Let $x=\sigma_{\text{std}}$ and $\displaystyle y=\bigwedge\nolimits^2\sigma_{\text{std}}$. Then $\text{ker }i^*=4x^3+4y^3+27-x^2y^2-18xy$. Note that $\text{dim }(4x^3+4y^3+27)=\text{dim }(x^2y^2+18xy)=243$. By viewing these two representations as unitary linear transformations on a 243-dimensional complex vector space, we can define $b:=\delta_S^{G/S}((4x^3+4y^3+27)(x^2y^2+18xy)^{-1})$, which generates the $R(S; \mathbb{Q})$-submodule $\mathcal{L}_S$. Following the last example, we also denote the maximal torus of diagonal matrices of $SU(3)$ by $T$. Then $(G/S)^S$ is $U/S$, where $\displaystyle U=\left.\left\{\begin{pmatrix} A&\\ &\det(A)^{-1}\end{pmatrix}\right|A\in U(2)\right\}$. Note that $U/S\cong PSU(2)$, which is homeomorphic to $\mathbb{RP}^3$. Let $a$ be an arbitrary element of $A$. The Atiyah-Segal localization formula for $\pi_*(ab)$ reads
	\[\pi_*(ab)=\pi_{U/S*}\left(\frac{i_{U/S, G/S}^*(ab)}{e_S(N_{U/S}^*)}\right).\]
	We will show that $i_{U/S, G/S}^*(b)=0$ and hence the index is $0\in I(S; \mathbb{Q})$. Consider the following commutative diagram.
	\[\xymatrix@+3pc{K_S^*(G/S)\ar[r]^{i_{U/S, G/S}^*}\ar[d]_{\pi_{G, G/S}^*}&K_S^*(U/S)\ar[r]^{i_{T/S, U/S}^*}\ar[d]_{\pi_{U, U/S}^*}& K_S^*(T/S)\ar[d]^{\pi_{T, T/S}^*}\\ K_S^*(G)\ar[r]^{i_{U, G}^*}&K_S^*(U)\ar[r]^{i_{T, U}^*}& K_S^*(T)\cong R(S)\otimes K^*(T)}\]
	Note that among the maps on the right square, $i_{T, U}^*$ and $\pi_{T, T/S}^*$ are injective while the kernel of $i_{T/S, U/S}^*$ and $\pi_{U, U/S}^*$ is generated by the 2-torsion element in the degree 0 piece represented by the reduced nontrivial complex line bundle over $\mathbb{RP}^3$. Note also that $a$ lives in the degree $-1$ piece of $K_S^*(G/S)$. Therefore to show that $i_{U/S, G/S}^*(b)=0$, it suffices to show that $i_{T, U}^*\circ i_{U, G}^*\circ \pi_{G, G/S}^*(a)=0$. We have
	\begin{align*}
		i_{T, U}^*\circ i_{U, G}^*\circ \pi_{G, G/S}^*(b)&=i_{T, U}^*\circ i_{U, G}^*\delta_S^G(4x^3+4y^3+27-x^2y^2-18xy)\\
		&=\delta_S^T(4(t_1+t_2+t_1^{-1}t_2^{-2})^3+4(t_1t_2+t_1^{-1}+t_2^{-1})^3+27\\
		&-(t_1+t_2+t_1^{-1}t_2^{-1})^2(t_1t_2+t_1^{-1}+t_2^{-1})^2-18(t_1+t_2+t_1t_2)(t_1t_2+t_1^{-1}t_2^{-1}))\\
		&(\text{Note that }i_{T, G}^*x=t_1+t_2+t_1^{-1}t_2^{-1}\text{ and }i_{T, G}^*y=t_1t_2+t_1^{-1}+t_2^{-1})\\
		&=\delta_S^T(-t_1^4t_2^2+2t_1^3t_2^3-t_1^2t_2^4+2t_1^3+2t_2^3-2t_1^2t_2-2t_1t_2^2-t_1^2t_2^{-2}-t_1^{-2}t_2^2\\
		&-2t_1t_2^{-1}-2t_1^{-1}t_2+6+2t_1^{-3}+2t_2^{-3}-2t_1^{-1}t_2^{-2}-2t_1^{-2}t_2^{-1}-t_1^{-2}t_2^{-4}-t_1^{-4}t_2^{-2}\\
		&+2t_1^{-3}t_2^{-3})
	\end{align*}
	Using $\delta_S^T(t_1^mt_2^n)=mt^{m+n}\otimes\delta(t_1)+nt^{m+n}\otimes\delta(t_2)\in R(S)\otimes K^*(T)$ (cf. \cite[Lemma 4.19]{F}), we have that the above expression is 0. Now the index $\pi_*(ab)$ is 0, which is in $I(S; \mathbb{Q})$. By Theorem \ref{indexthry}, we again obtain the conclusion that $(G, S)$ is not isotropy formal.
\end{example}
\begin{thebibliography}{BBFMR}
	\bibitem[ABC+]{ABC+}
	S. Agrawal, et al., 
	\emph{The CRing project: a collaborative, open source textbook on commutative algebra}, 2011. URL: \text{http://math.uchicago.edu/$\sim$amathew/CRing.pdf}.
	
	\bibitem[A]{A}
	M. F. Atiyah, 
	\emph{On the K-theory of compact Lie groups}, 
	Topology, Vol. 4, pp. 95--99, 1965.
	
	\bibitem[AM]{AM}
	M. F. Atiyah, I. G. Macdonald, 
	\emph{Introduction to commutative algebra}, 
	Addison-Wesley, 1969.
	
	\bibitem[BBFMP]{BBFMP}
	A. Borel, G. Bredon, E. E. Floyd, D. Montgomery, R. Palais, 
	\emph{Seminar on transformation groups}, Number 46 in Ann. of Math. Study Princeton University Press, 1960.	
	
	\bibitem[Bo]{Bo}
	N. Bourbaki, \emph{Lie groups and Lie algebras}, Chapters 7-9, Springer-Verlag, 2005.	
	
	\bibitem[BZ]{BZ}
	J.-L. Brylinski, B. Zhang, 
	\emph{Equivariant K-theory of compact connected Lie groups}, 
	K-theory, Vol. 20, pp. 23-36, 2000.
	
	\bibitem[Ca]{Ca}
	J. D. Carlson, 
	\emph{Equivariant formality of isotropic torus actions}, 
	J. Homotopy Relat. Struct. Vol. 14, pp.199--234, 2019. 
	
	\bibitem[CFok]{CFok} 
	J. D. Carlson, C.-K. Fok,  
	\emph{Equivariant formality of isotropy actions}, 
	J. London Math. Soc. (2), Vol. 97, no. 3, pp. 470--494, 2018.
	
	\bibitem[CH]{CH}
	J. D. Carlson, C. He, 
	\emph{Equivariant formality of corank-one isotropy actions and products of rational spheres}, 
	Math. Z. Vol. 305, no. 21, 2023.
	\bibitem[F]{F}
	C.-K. Fok, 
	\emph{The Real K-theory of compact Lie groups}, 
	Symmetry, integrability and geometry: Methods and Applications, Vol. 10, 022, 26 pages, 2014.
	
	\bibitem[F2]{F2}
C.-K. Fok, \emph{Equivariant formality in K-theory}, 
New York J. Math. Vol. 25, pp. 315--327, 2019.

	\bibitem[F3]{F3}
	C.-K. Fok, \emph{Cohomology and K-theory of compact Lie groups}, 
	Expository notes, available at \text{pi.math.cornell.edu/$\sim$ckfok/Cohomology\_Lie\_groups.pdf}.
	
	\bibitem[Go]{Go}
	O. Goertsches, \textit{The equivariant cohomology of isotropy actions on symmetric spaces}, Doc. Math., Vol. 17 (2012), pp. 79--94.
	
	\bibitem[GN]{GN} 
	O. Goertsches, S. H. Noshari, \textit{Equivariant formality of isotropy actions on homogeneous spaces defined by lie group automorphisms}, 
	J. Pure Appl. Algebra, Vol. 220, No. 5 (2016), pp. 2017--2028. 

	\bibitem[GKM]{GKM}
	M. Goresky, R. Kottwitz, R. MacPherson, 
	\emph{Equivariant cohomology, Koszul duality, and
the localization theorem}, Invent. Math., Vol. 131 (1998), pp. 25--83.

	\bibitem[GHV]{GHV}
	W. H. Greub, S. Halperin, and R. Vanstone,  
	\emph{Connections, curvature, and cohomology, vol. III:
Cohomology of principal bundles and homogeneous spaces}, Academic Press, 1976.	

	\bibitem[GD]{GD}
	A. Grothendieck, J. Dieudonn\'e, 
	\emph{El\'ements de g\'eom\'etrie alg\'ebrique}, 
	Publ. Math. IHES Vol. 20, pp. 116--153, 1964.
	

	\bibitem[H]{H}
	R. Hartshorne, 
	\emph{Algebraic Geometry}, 
	Graduate Texts in Mathematics 52, Springer, New York, 2013
			
	\bibitem[Ho]{Ho} 
	L. Hodgkin, 
	\emph{On the K-theory of Lie groups}, 
	Topology, Vol. 6, pp. 1--36, 1967. 
	
	 \bibitem[RK]{RK} 
	 I. Rosu, with an appendix by I.Rosu and A. Knutson, \textit{Equivariant K-theory and equivariant cohomology}, Math. Z., Vol. 243, pp. 423--448, 2008.
	 	
	\bibitem[Se]{Se}
	G. Segal, 
	\emph{The representation-ring of a compact Lie group}, 
	Publ. Math. Inst. Hautes \'Etudes Sci., Vol. 34, pp. 113--128, 1968.
	
	\bibitem[Se2]{Se2}
	G. Segal, 
	\emph{Equivariant K-theory}
	Publ. Math. Inst. Hautes \'Etudes Sci., Vol. 34, pp. 129-151, 1968.
	
	\bibitem[Sh]{Sh}
	H. Shiga, \textit{Equivariant de Rham cohomology of homogeneous spaces}, J. Pure Appl. Algebra, 106(2) (1996), pp. 173--183.
	
	\bibitem[ST]{ST}
	H. Shiga, H. Takahashi, \textit{Remarks on equivariant cohomology of homogeneous spaces}, Technical
report 17, Technological University of Nagaoka, May 1995.

	\bibitem[SP]{SP}
	The Stacks Project, \url{https://stacks.math.columbia.edu}, 2026.
		
	\end{thebibliography}
	\noindent\footnotesize{\textsc Xi'an Jiaotong-Liverpool University,\\
111 Ren’ai Road, Suzhou Industrial
Park,\\Suzhou, Jiangsu Province 215123, China}\\
\\
\textsc{E-mail}: \texttt{ChiKwong.Fok@xjtlu.edu.cn}\\
%\textsc{URL}: \texttt{https://sites.google.com/site/alexckfok}	
\end{document}
